% concatenated from main.tex
\documentclass[11pt,oneside,leqno]{amsart}
\usepackage[utf8]{inputenc}
\usepackage[english]{babel}
\usepackage{amssymb,amsmath,latexsym,amsthm,amsfonts,bbm,dsfont}
\usepackage{geometry}
\usepackage{float}
\usepackage{url}
\usepackage[colorlinks=true,citecolor=red,linkcolor=blue,pdfpagetransition=Blinds]{hyperref}
\numberwithin{equation}{section}
\usepackage{mathtools, stmaryrd}
\usepackage{natbib}
\usepackage{textgreek}
\usepackage{xcolor}
\numberwithin{equation}{section}
\usepackage{graphicx}
\newcommand{\taut}{\text{\texttau}}
\newcommand{\E}{\mathbb{E}}
\newcommand{\Or}{\mathcal{O}}
\usepackage{comment}
\newtheorem{theorem}{Theorem}[section]
\newtheorem{proposition}[theorem]{Proposition}
\newtheorem{lemma}[theorem]{Lemma}
\newtheorem{corollary}[theorem]{Corollary}
\newtheorem{remark}[theorem]{Remark}
\newtheorem{assumption}[theorem]{Assumption}
\newtheorem{example}[theorem]{Example}
%%%%%%%%%%%%%%%%%%%
\newcommand{\trp}{\mathtt t^+}
\newcommand{\trm}{\mathtt t^-}
\newcommand{\trpm}{\mathtt t^\pm}

%%%%%%%%%%%%%%%%%%%
%%%%%%%%%%%%%%%%%%%%%%%%5
\begin{document}
\title[Asymptotic behavior of Carleman weight functions]{Asymptotic behavior of Carleman weight functions and application to controllability}
\author{Ariel A. P\'erez}
\address{Departmento de Matem\'atica, Universidad del B\'io-B\'io, Avda. Collao 1202, Concepci\'on, Chile}
\email{aaperez@ubiobio.cl}
\subjclass[2020]{93B05,35R30,93B07,35B45}
\keywords{Carleman estimate, Control Theory, Inverse Problems, Asymptotic behavior}
\begin{abstract}
 In the development of controllability and inverse problem results for semi-discrete systems, by using Carleman estimates, it is required to estimate of the discrete operators applied to Carleman weight functions. This work aims to establish the asymptotic behavior of Carleman weight functions under these discrete operators. We provide a characterization of the error term in arbitrary order and dimension, extending previously known results. This generalization is of independent interest due to its applications in deriving Carleman estimates for semi-discrete stochastic operators. The aforementioned estimates hold for Carleman weight functions used for parabolic, hyperbolic, and elliptic operators, which are applied to obtain control and inverse problems results for those operators. We apply these results to obtain $\phi$-controllability result for a fully discrete parabolic operator, which is based in a Carleman estimate for a fully-discrete parabolic operator. 
\end{abstract}
\maketitle
%%%%%%%%%%%%%%%%%%%%
%-------------
\section{Introduction}
%-------------
 Let $d\geq 1$ and $T,L_{1}$, $\ldots$,  $L_{d}$ be positive real numbers; and $\Omega:=\Pi_{i=1}^{d}(0,L_{i})$, with $\omega\Subset \Omega$. Consider the following parabolic problem
\begin{equation}\label{system:control}
\begin{cases}
\partial_{t} y - \sum_{i=1}^{d}\partial_{x_{i}}( \gamma_{i}\partial_{x_{i}}y)  = \mathbbm{1}_\omega v, & (x,t) \in \Omega \times (0,T),\\
y(x,t) = 0, & (x,t) \in \partial\Omega \times (0,T),\\
y(x,0) = g(x), & x \in \Omega,
\end{cases}
\end{equation}
here $\gamma_{i}(x)>0$ for all $x\in\Omega$, and $\mathbbm{1}_\omega$ stands for the indicator function of the set $\omega$. The goal is to find $v\in L^{2}((0,T)\times\Omega)$ such that $y(T)=0$. In this case, it is said that the system \eqref{system:control} is null-controllable. The aforementioned controllability problem was solved independently by Lebeau and Robbiano \cite{lebeau1995controle} and Fursikov and Imanuvilov \cite{fursikov-1996}. The usual methodology is that the null-controllability result for the system \eqref{system:control} is equivalent, by a duality argument, to the following observability inequality
\begin{equation}\label{adjoint:continuous}
    |q(0)|^{2}_{L^{2}(\Omega)}\leq C_{\text{obs}}^{2}\left\| q\right\|^{2}_{L^{2}(\omega\times(0,T))},
\end{equation}
where $q$ is a solution to the adjoint system
\begin{equation}\label{system:control:adjoint}
\begin{cases}
\partial_{t} q + \sum_{i=1}^{d}\partial_{x_{i}}( \gamma_{i}\partial_{x_{i}}q)  = 0, & (x,t) \in \Omega \times (0,T),\\
q(x,t) = 0, & (x,t) \in \partial\Omega \times (0,T),\\
q(x,T)=q_{T}(x),&x\in \Omega.
\end{cases}
\end{equation}
A popular strategy to prove the inequality \eqref{adjoint:continuous} is based on a Carleman estimate for the system \eqref{system:control:adjoint}. The first application of these estimates was to prove unique continuation properties (see \cite{carleman1939}). Later, for control and inverse problems Carleman estimates have become a fundamental tool. A unified approach to Carleman estimates for second-order partial differential equations is the reference \cite{fu-2019} and references therein.
\par Unfortunately, it is not feasible to achieve the null-controllability for the discretization of system \eqref{system:control}. The next example shows that the null controllability generally does not hold for fully-discrete systems by building a non-observable function for a discretization of the adjoint system \eqref{system:control:adjoint}. To this end, let $M,N\in \mathbb{N}$, $T>0$, and define $h=1/(M+1)$ and $\Delta t=T/N$. Then, the uniform meshes in space and time are given by follows
\begin{align*}
    \mathcal{K}=\{x_i:=ih\,;\, i\in \{ 0,\ldots,M+1 \}\},\quad
    \mathcal{N}=\{t_j:=j\Delta t\, ;\, j\in \{ 0,\ldots,N-1 \}\}. 
\end{align*}
Thus, we consider a regular partition of the square $(0,1)^{2}$ defined by $(0,1)^{2}\cap\mathcal{M}^{2}$ where $\mathcal{M}:=(0,1)\cap\mathcal{K}$. Using the finite-difference scheme, the adjoint discrete system \eqref{system:control:adjoint} of the heat equation, for $d=2$, $\gamma_{i}(x)=1$ and $L_{i}=1
$, can be written as
\begin{equation}\label{sys:discrete:adjoint}
\begin{split}
    D_{t}q_{i,j}+\Delta_{h}\trp q_{i,j}&=0 \quad \forall(x_{i},x_{j},t_{j})\in\mathcal{M}^{2}\times\mathcal{N},\\
    q_{i,0}=q_{i,M+1}&=0\  \ \ \forall (x_{i},t_{j})\in\mathcal{M}\times \mathcal{N},\\
    q_{0,j}=q_{M+1,j}&=0 \ \ \ \forall (x_{j},t_{j})\in \mathcal{M}\times\mathcal{N},\\
    \end{split}
\end{equation}
where the five-point approximation of the Laplacian operator in uniform meshes is utilized, as given by
\begin{equation}
    \Delta_{h}q_{i,j}:=\frac{1}{h^{2}}\left(q_{i+1,j}+q_{i,j+1}+q_{i-1,j}+q_{i,j-1}-4q_{i,j}\right),
\end{equation}
and $D_{t} q:=\dfrac{\trm q-\trm q}{\Delta t}$ in $\mathcal{N}$, where $\trpm q(t):=q(t\pm \Delta t/2)$ and $q_{ij}=q(x_{i},x_{j})$. Combining the ideas of \cite{zuazua:2005} and \cite{ervedoza:2019}, we can build a non-observable solution of the system \eqref{sys:discrete:adjoint}. 
Indeed, notice that for
\begin{equation}
    \tilde{q}_{i,j}:=\begin{cases}
    \ \ 1,& i=j \text{ even,}\\
    \ \ 0,& i\ne j,\\
    -1,& i=j\text{ odd},\\
    \end{cases}
\end{equation}
it follows that 
\begin{equation}
    \Delta_{h}\tilde{q}_{i,j}=\begin{cases}
    -\frac{4}{h^{2}},& i=j \text{ even,}\\
    \ \ \  0,& i\ne j, \\
    \ \ \frac{4}{h^{2}},& i=j\text{ odd,}.\\
    \end{cases}=-\frac{4}{h^{2}}\tilde{q}_{i,j}.
\end{equation}
Then $\tilde{q}_{i,j}$ is an eigenfunction of $\Delta_{h}$, with eigenvalue $-4/h^{2}$. In turn, for $\alpha \Delta t<1$ and with 
$\displaystyle \mathfrak{a}:=\left( \frac{1}{1-\alpha\Delta t}\right)^{1/\alpha\Delta t}$,  it obtained that $D_{t}\mathfrak{a}^{\alpha t}=\alpha\trp \mathfrak{a}^{\alpha t}$. Thus, 
considering $\alpha=4/h^{2}$ such that $\alpha \Delta t<1$ it follows that $\overline{q}_{i,j}:=\mathfrak{a}^{-4(T-t)/h^{2}}\tilde{q}_{i,j}$ solves \eqref{sys:discrete:adjoint} since $D_{t}\overline{q}=\frac{4}{h^{2}}\trp\overline{q}_{i,j}$. Therefore, if any diagonal nodes do not belong to $\omega$ yields $\left\| \overline{q}_{i,j}\right\|_{L^{2}_{h}(\omega\times (0,T))}=0$ although the norm of $\overline{q}(x,0)$ has size $\mathfrak{a}^{-C/h^{2}}$. This implies that a discrete observability inequality like the continuous setting \eqref{adjoint:continuous}, cannot be achieved for a fully-discrete system \eqref{sys:discrete:adjoint}.
\par Due to the aforementioned counter-example, achieving null controllability in the fully or semi-discrete setting is not possible in general, even though an alternative objective is to build bounded controls such that at a given time the norm of the solution of the system decreases exponentially, as the discrete parameters $\Delta t$ and/or $h$ tend towards zero. That is, in the case of space semi-discrete system, to build semi-discrete control $u$ such that at time $T$ the solution of the system $y$ (with an appropriate norm on the space of solutions $X$) verifies
\begin{equation}
    \left\| y(x,T)\right\|_{X} \approx C\sqrt{\phi(h)},
\end{equation}
with $C>0$, $\left\| u\right\|_{Y}\leq C$ uniformly respect to $h$ (with the appropriate norm in the space of controls $Y$), and
where $h\longrightarrow \phi(h)$ is any given function of the discretization parameter such that
\begin{equation}\label{eq:phi}
    \liminf_{h\rightarrow 0}\frac{\phi(h)}{e^{-c/h}}>0.
\end{equation}
This approach is known as $\phi$-null controllability for the space semi-discrete system, and also can be formulated for time-discrete systems and/or fully-discrete systems.
 
\par In several setting, the $\phi$-null controllability has been studied and the analysis is based on the penalization Hilbert uniqueness method, we refer to the work \cite{boyer-canum} for an overview of the method for space semi-discrete parabolic setting. In particular, the works \cite{boyer-2014,boyer-2010-1d-elliptic,BHLR:2010} establish $\phi$-null controllability results for semi-discrete parabolic operators in the spatial semi-discrete setting. For discontinuous diffusion coefficient we refer to \cite{Thuy} and to \cite{Thuy2} for the setting of space semi-discrete parabolic operator in Banach space. Also in the space-discrete setting in \cite{BHSDT:2019}, the authors examined insensitizing controls for a semilinear parabolic equation in one dimension. For higher order operator we refer to \cite{CLTP-2022} and \cite{kumar:hal-04828698}, where the boundary $\phi$-null controllability is studied. Likewise, in  \cite{BHS:2020} is presented the $\varphi(\Delta t)$-null controllability within the time-discrete framework for the time-discrete heat equation and also for some parabolic systems. The same kind of controllability for a simplified stabilized Kuramoto-Sivashinsky system was also addressed in \cite{HS:2023} and recently for parabolic system with Kirchhoff boundary condition in \cite{bhandari:hal-04597449}. As for fully-discrete systems, an important contribution was made in \cite{GC-HS-2021} for null Dirichlet boundary conditions and in \cite{LMPZ:2023} for dynamic boundary conditions both of them in the one-dimensional setting.

\par In all the aforementioned works, their results follow from suitable Carleman estimates, and they need an estimate of the discrete difference and average operator applied on the Carleman weight functions. It is worth noticing that, in arbitrary dimension, there are no results for general computation of the discrete operators over the continuous Carleman weight functions; see \cite{LOP-2020} for the one-dimensional case. Hence, in this work, we aim to cover that gap by giving general results about the asymptotic behavior of the Carleman weight functions; see Proposition \ref{prop:weight}, Theorems \ref{theo:weight:estimates}, \ref{theo:time:derivate} and \ref{theo:fully:weight:estimates}. We emphasize that our introduction is focusing on controllability issues but our scope also reaches semi-discrete inverse problem and unique continuation properties for hyperbolic and elliptic operators respectively; see Section \ref{sec:concluding} for some remarks in these directions. \\
\par Almost all the fully-discrete $\phi(h)$-controllability results are in 1-D setting. The main difficulty to extend these results into higher-dimension is that there is no general asymptotic behavior of Carleman weight function. Then, adopting the same notation from \cite{LDOP-2021,LPP:2026,LPP:2026b,LLRP:2025} for the sets $\mathcal{M}$, $\mathcal{M}_{i}^{\ast}$, $\partial_{i}\mathcal{M}$, $\mathcal{N}$, the operators $D_{i}$, $A_{i}$, $t_{r}^{i}$, we show how to use Theorems \ref{theo:weight:estimates}, \ref{theo:time:derivate} and \ref{theo:fully:weight:estimates} to obtain a fully discrete Carleman estimate which implies $\phi-$controllability result for the following system
\begin{equation}\label{eq:controlled}
\begin{cases}
D_{t}y-\sum_{i=1}^{d}D_{i}(\gamma_{i}D_{i}\trp y)+\trp(a\,y)=1_{\omega}v,&\forall(x,t)\in\mathcal{M}\times\mathcal{N}^{\ast},\\
y\big|_{\partial\mathcal{M}}=0,&\forall t\in\mathcal{N}^{\ast},\\
y(x,0)=y_{0}(x),&\forall x\in\mathcal{M},
\end{cases}
\end{equation}
where  $y\in C(\overline{\mathcal{M}}\times\overline{\mathcal{N}})$, $y_{0}\in C(\mathcal{M})$ is the initial datum, $v\in C(\omega\times\mathcal{N}^{\ast})$ is the control function acting on a discrete subset $\omega\subset\mathcal{M}$, $a=a(x,t)$ is a bounded potential, and $\Gamma(x):= \mathrm{Diag}(\gamma_{1}(x),\gamma_{2}(x),\ldots, \gamma_{d}(x))$, $\gamma_i>0$ for all $i=1,\ldots,d$ and 
$$
 \mbox{reg}(\gamma):= \text{ess}\sup_{\substack{x\in \mathcal{M} \\ i=1,\ldots,d}}\left(\gamma_i+\frac{1}{\gamma_i}+{ \sum_{j=1}^{d}|\partial_{i}\gamma_i|^2}\right)<c_0.
$$
for some constant $c_0> 0$. Indeed, following the penalized HUM strategy, see \cite{boyer-canum}, we have the $\phi$-controllability result for the system \eqref{eq:controlled}
\begin{theorem}\label{thm:phi:controllability}
Let $T>0$, $\mu\geq 1$, and $h$, $\Delta t$ as in Theorem~\ref{thm:observability}. Then, for any initial datum $y_{0}\in C(\mathcal{M})$, there exists a control function $v$ satisfying
\begin{equation}\label{eq:control:bound}
\|v\|_{L^{2}_{h}(\omega\times\mathcal{N}^{\ast})}\leq\sqrt{C_{\text{obs}}}\,\|y_{0}\|_{L^{2}_{h}(\mathcal{M})},
\end{equation}
and the solution of the controlled system \eqref{eq:controlled} satisfies
\begin{equation}\label{eq:phi:null}
\|y(T)\|_{L^{2}_{h}(\mathcal{M})}\leq\sqrt{C_{\text{obs}}}\,\sqrt{\phi(h)}\,\|y_{0}\|_{L^{2}_{h}(\mathcal{M})},\qquad\phi(h)=e^{-C_{1}/h^{\min\{\mu/4,1\}}},
\end{equation}
provided $h\leq\min\{h_{0},h_{1}\}$ and $\Delta t\leq\min\{T^{-2}h^{\mu},(4\|a\|_{L^{\infty}_{h}})^{-1}\}$, \\with $C_{\mathrm{obs}}=e^{C(1+1/T+\|a\|_{L^{\infty}_{h}}^{2/3}+T\|a\|_{L^{\infty}_{h}})}$.
\end{theorem}
In \cite[section 4]{LMPZ:2023} is the strategy to follow to prove Theorem \ref{thm:phi:controllability}. Moreover, Theorem \ref{thm:phi:controllability} can be extended to more general functions $\phi$ satisfying
\begin{equation*}
\liminf_{h\to 0}\phi(h)\,e^{C_{1}/h^{\min\{\mu/4,1\}}}>0.
\end{equation*}
Indeed, for such $\phi$, there exists $h^{\ast}>0$ such that $e^{-C_{1}/h^{m_{\mu}}}\leq\phi(h)$ for all $h\leq h^{\ast}$. Then, Theorem~\ref{thm:observability} holds with $\phi(h)$ replacing $e^{-C_{1}/h^{m_{\mu}}}$, and the same penalised HUM argument applies. Here, Theorem \ref{thm:observability} is the observability inequality of  the corresponding adjoint system is given by
\begin{equation}\label{eq:adjoint}
\begin{cases}
D_{t}q+\sum_{i=1}^{d}D_{i}(\gamma_{i}D_{i}\trm q)+\trm(a\,q)=0,&\forall(x,t)\in\mathcal{M}\times\mathcal{N},\\
q\big|_{\partial\mathcal{M}}=0,&\forall t\in\mathcal{N},\\
q(x,T+\Delta t/2)=q_{T}(x),&\forall x\in\mathcal{M}.
\end{cases}
\end{equation}
 for $q\in C(\overline{\mathcal{M}}\times\overline{\mathcal{N}}^{\ast})$.\\
 
 
\par The organization of this work is as follows. In section \ref{preliminary} we provided some results concerning the average and difference operators and the asymptotic behavior of smooth functions. In section \ref{sec:asymp:weight} we apply the results from the previous section into the Carleman weight function. In section \ref{sec:Carleman} we prove the Carleman estimate for the system \ref{eq:adjoint} although several computation are in section \ref{sec:proof:carleman}. Section \ref{sec:concluding} provides concluding remarks about the potential applicability of the results presented in section \ref{sec:asymp:weight} for deterministic and stochastic operators.
%%%%%%%%%%%%%%%%%%%
\section{Asymptotic behavior of smooth functions}\label{preliminary}
In this section, we aim to establish the asymptotic behavior of a smooth function which will be applied in particular in the next section to Carleman weight functions. To this end, we consider functions continuously defined in a neighborhood of $\overline{\Omega}$. We define the average and difference operators, in the direction $e_{i}$, with $\{e_{i}\}_{i=1}^{d}$ being
the usual basis of $\mathbb{R}^{d}$, as 
\begin{equation}
\begin{split}
    A_{i}u(x)&:=\frac{1}{2}\left( u(x+\frac{h}{2}e_{i})+u(x-\frac{h}{2}e_{i})\right),\\
    D_{i}u(x)&:=\frac{1}{h}\left(u(x+\frac{h}{2}e_{i})-u(x-\frac{h}{2}e_{i}) \right),
    \end{split}
\end{equation}
  For these operators, we have the following product rule.
\begin{lemma}[{\cite[Lemma 2.1 and 2.2]{BHLR:2010}}] \label{base:case}
    Let the functions $u$ and $v$ be continuously defined in a neighborhood of $\overline{\Omega}$. For $i\in\{1,\ldots,d\}$, we have the difference and average operators satisfy
    \begin{align}
        D_{i}(uv)=&D_{i}u\,A_{i}v+D_{i}v\,A_{i}u,\label{D:product}\\
        A_{i}(uv)=&A_{i}u\,A_{i}v+\frac{h^{2}}{4}D_{i}u\,D_{i}v.\label{A:product}
    \end{align}
\end{lemma}
%------------------
It is possible to generalize the above product rule for the difference and average operators.
\begin{proposition}
Let $u$ and $v$ be two smooth functions and $n\in\mathbb{N}$. Then, it holds a Leibniz rule for the difference operator given by
\begin{equation}\label{eq:discreteleibniz}
    D_{i}^{n}(u\,v)=\sum_{k=0}^{n}\binom{n}{k}D_{i}^{n-k}A_{i}^{k}u\,A^{n-k}_{i}D^{k}_{i}v.
\end{equation}
Moreover, for the average operator, we have
\begin{equation}\label{eq:discrete:leibniz:average}
    A_{i}^{2m}(uv)=\sum_{j=0}^{m}\frac{h^{2j}}{2^{2j}}D_{i}^{2j}(uv),\quad\forall m\in\mathbb{N}.
\end{equation}
\end{proposition}
\begin{proof}
    The proof follows by induction, with the base case given by Lemma \ref{base:case}.
\end{proof}
\begin{remark}
    Note that \eqref{eq:discrete:leibniz:average} is stated for an even power, although by using \eqref{A:product} the odd case can also be handled.
\end{remark}
%-----------------
Our first result states the asymptotic behavior of a smooth function when several iterations of the discrete operator are applied, which is a  higher-dimension extension of \cite[Proposition 4.1]{LOP-2020}.
\begin{proposition}\label{prop:difference:average}
Let $f$ be a smooth function defined in a neighborhood of $\overline{\Omega}$. For 
$i\in\{ 1,\ldots,d\}$ we have
\begin{equation}\label{eq:ave:diff}
    \begin{split}
        A^{n}_{i}f=&f+R_{A_{i}^{n}}(f),\quad \forall n\in \mathbb{N},\\
        D^{n}_{i}f=&\partial_{i}^{n}f
        +R_{D_{i}^{n}}(f),\quad\forall n\in\mathbb{N
        }.
    \end{split}
\end{equation}
where  $\displaystyle R_{D_{i}^{n}}(f):=h^{2}\sum_{k=0}^{n}\binom{n}{k}(-1)^{k}\left(\frac{n-2k}{2} \right)^{n+2}\int_{0}^{1}\frac{(1-\sigma)^{n+1}}{(n+1)!}\partial_{i}^{n+2}f(\cdot+\tilde{e}_{i}\sigma)d\sigma$ and \\ $\displaystyle R_{A_{i}^{n}}(f):=\frac{h^{2}}{2^{n}}\sum_{k=0}^{n}\binom{n}{k}\frac{(n-2k)^{2}}{4}\int_{0}^{1}(1-\sigma)\partial_{i}^{2}f(\cdot+\tilde{e}_{i}\sigma)d\sigma$,
with $\tilde{e}_{i}:=\frac{(n-2k)h}{2}e_{i}$ and $\sigma\in[0,1]$.
\end{proposition}
\begin{proof}
The proof follows a similar strategy developed in the proof of \cite[Proposition 4.1]{LOP-2020} with the Taylor expansion \begin{equation}\label{eq:taylor}
 f(x+\eta)=\sum_{j=0}^{l-1}\frac{1}{j!}f^{(j)}(x;\eta,\ldots, \eta)+\int_{0}^{1}\frac{(1-\sigma)^{l-1}}{(l-1)!}f^{(l)}(x+\sigma \eta;\eta,\ldots,\eta)d \sigma.
\end{equation} Indeed, denoting $\displaystyle\tau_{(n-2k)i}f(x):=f\left(x+\frac{(n-2k)h}{2}e_{i}\right)$, we use \eqref{eq:taylor} with $l=2$, $\displaystyle\eta=\frac{(n-2k)h}{2}e_{i}$ to obtain
\begin{equation*}
    \tau_{(n-2k)i}f=f+\frac{(n-2k)h}{2}\partial_{i}f+\frac{(n-2k)^{2}h^{2}}{4}\int_{0}^{1}(1-\sigma)\partial_{i}^{2}f(\cdot+\frac{(n-2k)h}{2}e_{i}\sigma)d\sigma.
\end{equation*}
Then, using the above identity for the average operator, we have 
\begin{align*}
    A_{i}^{n}f=&\frac{1}{2^{n}}\left(\tau_{i}+\tau_{-i} \right) ^{n}f\\
        =&\frac{1}{2^{n}}\sum_{k=0}^{n}\binom{n}{k}\tau_{(n-2k)i}f\\
        =&\frac{1}{2^{n}}\sum_{k=0}^{n}\binom{n}{k}\left(f+\frac{(n-2k)h}{2}\partial_{i}f\right)+R_{A_{i}^{n}}(f)
\end{align*}
Thus, using that $\displaystyle\sum_{k=0}^{n}\binom{n}{k}=2^{n}$ and $\displaystyle\sum_{k=0}^{n}\binom{n}{k}k=n2^{n-1}$ we conclude that
\begin{equation*}
    A_{i}^{n}f=f+R_{A_{i}^{n}}(f),
\end{equation*}
which is the first identity from \eqref{eq:ave:diff}.\\    In turn, for the second identity from \eqref{eq:ave:diff}, we proceed similarly. Using \eqref{eq:taylor}, with $l=n+2$ and $\displaystyle\eta=\frac{(n-2k)h}{2}e_{i}$, it follows that
\begin{align}\label{eq:Df}
    D^{n}_{i}f=&\frac{1}{h^{n}}\left(\tau_{i}-\tau_{-i} \right)^{n}f\notag\\
        =&\frac{1}{h^{n}}\sum_{k=0}^{n}\binom{n}{k}(-1)^{k}\tau_{(n-2k)i}f\\
        =&\frac{1}{h^{n}}\sum_{k=0}^{n}\binom{n}{k}(-1)^{k}\sum_{j=0}^{n+1}\frac{1}{j!}\left(\frac{(n-2k)h}{2} \right)^{j}\partial_{i}^{j}f+R_{D_{i}^{n}}(f).\notag
        \end{align}
Now, using from \cite[Equation 2.2]{egorychev} the Tepper identity  
$$\displaystyle\sum_{k=0}^{n}(-1)^{k}\binom{n}{k}(x-k)^{r}=\begin{cases}
    0,&\ 0\leq r<n, \\
    n!,&\ r=n, 
\end{cases}$$ 
we notice that the non-zero terms, in the latter expression of \eqref{eq:Df}, are only for $j=n$ and $j=n+1$. Moreover, we have 
  \begin{equation}\label{eq:tepper:n+1}
      \mathcal{S}:=\sum_{k=0}^{n}(-1)^{k}\binom{n}{k}\left(\frac{n}{2}-k\right)^{n+1}=0.   
  \end{equation}
  Indeed, notice that applying the substitution $k\mapsto n-k$, we obtain
  \begin{equation*}
      \mathcal{S}=\sum_{k=0}^{n}(-1)^{n-k}\binom{n}{n-k}\left(\frac{n}{2}-(n-k)\right)^{n+1}=(-1)^{n}(-1)^{n+1}\sum_{k=0}^{n}(-1)^{k}\binom{n}{k}\left(\frac{n}{2}-k\right)^{n+1}=-\mathcal{S}
  \end{equation*}
  Hence, \eqref{eq:tepper:n+1} holds.
  Thus, by using \eqref{eq:tepper:n+1} it follows that \eqref{eq:Df} becomes
        \begin{equation}
        \begin{split}
        D_{i}^{n}f=&\frac{1}{h^{n}}\sum_{k=0}^{n}\binom{n}{k}(-1)^{k}\sum_{j=0}^{n+1}\frac{1}{j!}\left(\frac{(n-2k)h}{2} \right)^{j}\partial_{i}^{j}f+R_{D_{i}^{n}}(f)\\
        =&\frac{1}{h^{n}}\sum_{j=0}^{n+1}\frac{1}{j!}h^{j}\sum_{k=0}^{n}\binom{n}{k}(-1)^{k}\left(\frac{n}{2}-k \right)^{j}\partial_{i}^{j}f+R_{D_{i}^{n}}(f)\\
        =&\frac{1}{h^{n}n!}h^{n}\sum_{k=0}^{n}\binom{n}{k}(-1)^{k}\left(\frac{n}{2}-k \right)^{n}\partial_{i}^{n}f\\
        &+\frac{1}{h^{n}}\frac{1}{(n+1)!}h^{n+1}\sum_{k=0}^{n}\binom{n}{k}(-1)^{k}\left(\frac{n}{2}-k \right)^{n+1}\partial_{i}^{n+1}f+R_{D_{i}^{n}}(f)\\
        =&\partial_{i}^{n}f+R_{D_{i}^{n}}(f),
    \end{split}
\end{equation} 
which completes the proof.
\end{proof}
A combination of the estimates from \eqref{eq:ave:diff} yields estimates when the composition of the discrete operators is made in different directions.
\begin{corollary}\label{cor:difference:average}
Let $f$ be a smooth function defined in a neighborhood of $\overline{\Omega}$. For 
$i\ne j\in\{ 1,\ldots,d\}$ and $m,n\in\mathbb{N}$ we have 
\begin{equation}
    \begin{split}
        A_{j}^{m}A_{i}^{n}f=&f+R_{A_{j}^{m}}(f)+R_{A_{i}^{n}}(f)+R_{A_{j}^{m}A_{i}^{n}}(f),\\
        D_{j}^{m}D_{i}^{n}f=&\partial_{j}^{m}\partial_{i}^{n}f+R_{D_{j}^{m}}(f)+R_{D_{i}^{n}}(f)+R_{D_{j}^{m}D_{i}^{n}}(f),
    \end{split}
\end{equation}
and 
\begin{equation}
D_{i}^{n}A_{j}^{m}f=\partial_{i}^{n}f+R_{D_{i}^{n}}(f)+R_{A_{j}^{m}}(\partial^{n}_{i}f)+R_{A_{j}^{m}D_{i}^{n}}(f);
\end{equation}
where
\begin{equation*}
    R_{A^{m}_{j}A^{n}_{i}}(f):=\sum_{k=0}^{n}\sum_{k'=0}^{m}a_{k,k'}\int_{0}^{1}\int_{0}^{1}(1-\sigma)(1-\sigma^{\prime})\partial_{j}^{2}\partial_{i}^{2}f\left(\cdot+\tilde{e}_{ij}\right)d\sigma'd\sigma,
\end{equation*}
\begin{equation*}
    R_{D^{m}_{j}D^{n}_{i}}(f):=\sum_{k=0}^{n}\sum_{k'=0}^{m}b_{k,k'}\int_{0}^{1}\int_{0}^{1}(1-\sigma)^{n+1}(1-\sigma')^{m+1}\partial_{j}^{m+2}\partial_{i}^{n+2}f\left(\cdot+\tilde{e}_{ij}\right)d\sigma'd\sigma,
\end{equation*}
 and
 \begin{equation*}
    R_{D^{n}_{i}A^{m}_{j}}(f):=\sum_{k=0}^{n}\sum_{k'=0}^{m}c_{k,k'}\int_{0}^{1}\int_{0}^{1}(1-\sigma)^{2}(1-\sigma')^{n+2}\partial_{i}^{n+2}\partial_{j}^{2}f\left(\cdot+\tilde{e}_{ij}\right)d\sigma'd\sigma;
\end{equation*}
with $\tilde{e}_{ij}=\tilde{e}_{ij}(n,k,h,\sigma,\sigma'):=(n-2k)\frac{h}{2}e_{i}\sigma+(m-2k')\frac{h}{2}e_{j}\sigma'$ for $\sigma,\sigma'\in[0,1]$ and  $$a_{k,k'}:=\frac{h^{4}}{2^{m+n}}\binom{m}{k'}\binom{n}{k}\frac{(n-2k)^{2}}{4}\frac{(m-2k')^{2}}{4},$$
$$b_{k,k'}:=\frac{(-1)^{k+k'}h^{4}}{(n+1)!(m+1)!}\binom{m}{k'}\binom{n}{k}\left(\frac{n-2k}{2}\right)^{n+2}\left(\frac{m-2k'}{2}\right)^{m+2},$$
and 
$$c_{k,k'}:=\frac{(-1)^{k}h^{4}}{(n+1)!2^{m}}\binom{m}{k'}\binom{n}{k}\left(\frac{n-2k}{2}\right)^{n+2}\left(\frac{m-2k'}{2}\right)^{2}.$$
\end{corollary}
%%%%%%%%%%%%%%%%%%%%%
\section{Asymptotic behavior of Carleman weight functions}\label{sec:asymp:weight}
This section is devoted to extending the results on discrete operations performed on the Carleman weight functions, presented in \cite{LOP-2020}, to arbitrary dimensions. We set the weight function of the form $r = e^{s\varphi}$ and $\rho = r^{-1}$, with $s\geq 1$ a parameter; where    $\varphi(x)=e^{\lambda\psi(x)}$ and the function $\psi$ fulfills suitable assumptions concerning the specific issue; that is, controllability or inverse problems. Our main objective is to obtain an asymptotic behavior describing the dependence on $s, h$ in the following basic estimates in domain $\Omega$ when the discrete operators are applied. 
Let us recall that a multi-index has the form $\alpha:=(\alpha_{1},\ldots,\alpha_{n})\in\mathbb{N}^{n}$. Then, for $\alpha\in\mathbb{N}^{n}$ and $\xi\in\mathbb{R}^{n}$ we write the order of the multi-index as $|\alpha|:=\alpha_{1}+\cdots+\alpha_{n}$, the compact notation of a higher-order operator as $
\partial^{\alpha}:=\partial_{x_{1}}^{\alpha_{1}}\cdots\partial_{x_{n}}^{\alpha_{n}}$ and we also adopt the notation
$\xi^{\alpha}:=\xi_{1}^{\alpha_{1}}\cdots\xi_{n}^{\alpha_{n}}$.

Now, we consider two fundamental estimates for our weight function. 
\begin{lemma}[{\cite[Lemma 3.7]{BHLR:2010}}] \label{lem:derivative:wrt:x}
Let $\alpha$ and $\beta$ multi-indices. We have
\begin{align*}
\partial^{\beta}(r\partial^{\alpha} \rho)=&|\alpha|^{|\beta|}(-s\varphi)^{|\alpha|}\lambda^{|\alpha+\beta|}(\nabla \psi)^{\alpha+\beta}+|\alpha||\beta|(s\varphi)^{|\alpha|}\lambda^{|\alpha+\beta|-1}\mathcal{O}(1)\\
&+s^{|\alpha|-1}|\alpha|(|\alpha|-1)\mathcal{O}_{\lambda}(1)\\
=&\mathcal{O}_{\lambda}(s^{|\alpha|}).
\end{align*}
Let $\sigma\in [-1,1]$ and $i\in\{1,\ldots,d\}$. We have
\begin{equation}
\begin{split}
\partial^{\beta}\left( r(x)(\partial^{\alpha}\rho)(x+\sigma he_{i})\right)=\mathcal{O}_{\lambda}(s^{|\alpha|}(1+(sh)^{|\beta|}))e^{\mathcal{O}_{\lambda}(sh)}.
\end{split}
\end{equation}
Provided $sh\leq1$ we have $\partial^{\beta}\left( r(x)(\partial^{\alpha}\rho)(x+\sigma he_{i})\right)=\mathcal{O}_{\lambda}(s^{|\alpha|})$. The same expressions hold with $r$ and $\rho$ interchanged and with $s$ changed to $-s$.
\end{lemma}
Another useful estimate is the following. 
\begin{corollary}[{\cite[Corollary 3.8]{BHLR:2010}}]\label{Cor:derivative:wrt:x:2r}
Let $\alpha,\beta$ and $\delta$ multi-indices. We have
\begin{equation}
\begin{split}
\partial^{\delta}\left(r^{2}\,\partial^{\alpha}\rho\, \partial^{\beta}\rho\right)=&|\alpha+\beta|^{|\delta|}(-s\varphi)^{|\alpha+\beta|}\lambda^{|\alpha+\beta+\delta|}(\nabla\psi)^{\alpha+\beta+\delta}\\
&+|\delta||\alpha+\beta|(s\varphi)^{|\alpha+\beta|}\lambda^{|\alpha+\beta+\delta|-1}\mathcal{O}(1)\\
&+s^{|\alpha+\beta|-1}\left( |\alpha|(|\alpha|-1)+|\beta|(|\beta|-1)\right)\mathcal{O}(1)\\
&=\mathcal{O}_{\lambda}(s^{|\alpha+\beta|}).
\end{split}
\end{equation}
\end{corollary}
The proofs of the above results can be found in \cite{BHLR:2010}. We aim to obtain a general estimate for the discrete operators applied to the Carleman weight function.
 Mimicking the notation for multi-indices in the continuous setting, we define for the two-dimensional index $k=(k_{i},k_{j})$, with $k_{i},k_{j}\in\mathbb{N}$, the discrete operator $\mathbf{D}_{h}^{k}:=D_{i}^{k_{i}}D_{j}^{k_{j}}$. Analogously, for the average operator, we define $\mathbf{A}_{h}^{l}:=A_{i}^{l_{i}}A_{j}^{l_{j}}$ with $l:=(l_{i},l_{j})$ being a two-dimensional multi-index. Considering the previous notation and Proposition \ref{prop:difference:average}, we can extend \cite[Proposition 4.5]{LOP-2020} to arbitrary dimensions.
\begin{proposition}\label{prop:weight}
    Let $\alpha$ be a multi-index and $k,l$ be two-dimensional multi-indices. Let $i,j\in\{1,\ldots,d\} $; provided that $sh\leq 1$, we have
    \begin{equation}\label{eq:average:difference:weight}
\begin{aligned}
r\mathbf{D}_{h}^{k}\mathbf{A}_{h}^{l}\partial^{\alpha}\rho &= r\partial^{k+\alpha}\rho + s^{|\alpha|+k_{i}}\mathcal{O}_{\lambda}((sh)^{2}) \\
&\quad + s^{|\alpha|+k_{j}}\mathcal{O}_{\lambda}((sh)^{2}) + s^{|\alpha+k|}\mathcal{O}_{\lambda}((sh)^{2}).
\end{aligned}
\end{equation}
\end{proposition}
\begin{proof}
Applying Corollary \ref{cor:difference:average} to $\mathbf{A}_{h}^{l}\partial^{\alpha}\rho$, we write
\begin{equation}\label{eq:estimation:DA}
\begin{aligned}
r\mathbf{D}_{h}^{k}\mathbf{A}_{h}^{l}\partial^{\alpha}\rho=& r\partial^{k}(\mathbf{A}_{h}^{l}\partial^{\alpha}\rho)+rR_{D^{k_{j}}_{j}}(\mathbf{A}_{h}^{l}\partial^{\alpha}\rho)+rR_{D^{k_{i
}}_{i}}(\mathbf{A}_{h}^{l}\partial^{\alpha}\rho)\\
&+rR_{\mathbf{D}_{h}^{k}}(\mathbf{A}_{h}^{l}\partial^{\alpha}\rho).
\end{aligned}
\end{equation}
Then, using the first identity of Corollary \ref{cor:difference:average}, it follows for the first term on the right-hand side of \eqref{eq:estimation:DA} that
\begin{align*}
r\partial^{k}(\mathbf{A}_{h}^{l}\partial^{\alpha}\rho)=&r\partial^{k+\alpha}\rho+rR_{A_{j}^{l_{j}}}(\partial^{\alpha+k}\rho)+rR_{A_{i}^{l_{i}}}(\partial^{\alpha+k}\rho)+rR_{\mathbf{A}_{h}^{l}}(\partial^{\alpha+k}\rho)
\end{align*}
Now, our next task is to estimate the last three terms on the right-hand side above. Using Corollary \ref{cor:difference:average} and Lemma \ref{lem:derivative:wrt:x} we obtain the estimates
\begin{align*}
    rR_{A_{j}^{l_{j}}}(\partial^{\alpha+k}\rho)=&s^{|\alpha|+|k|}\mathcal{O}_{\lambda}((sh)^{2}),\\
    rR_{A_{i}^{l_{i}}}(\partial^{\alpha+k}\rho)=&s^{|\alpha|+|k|}\mathcal{O}_{\lambda}((sh)^{2}),\\
    rR_{\mathbf{A}_{h}^{l}}(\partial^{\alpha+k}\rho)=&s^{|\alpha|+|k|}\mathcal{O}_{\lambda}((sh)^{4}),
\end{align*}
which  yields 
\begin{align}\label{eq:estimation}
    r\partial^{k}(\mathbf{A}_{h}^{l}\partial^{\alpha}\rho)=r\partial^{k+\alpha}\rho+s^{|\alpha|+|k|}\mathcal{O}_{\lambda}((sh)^{2}).
\end{align} 
Similarly, we have the estimates 
\begin{equation}\label{eq:estimations}
\begin{aligned}
    rR_{D^{k_{j}}_{j}}(\mathbf{A}_{h}^{l}\partial^{\alpha}\rho)=&s^{k_{j}+|\alpha|}\mathcal{O}_{\lambda}((sh)^{2}),\\
    rR_{D^{k_{i}}_{i}}(\mathbf{A}_{h}^{l}\partial^{\alpha}\rho)=&s^{k_{i}+|\alpha|}\mathcal{O}_{\lambda}((sh)^{2}),\\
    rR_{\mathbf{D}_{h}^{k}}(\mathbf{A}_{h}^{l}\partial^{\alpha}\rho)=&s^{|k|+|\alpha|}\mathcal{O}_{\lambda}((sh)^{4}).
\end{aligned}
\end{equation}
Thus, combining the estimates \eqref{eq:estimations} and \eqref{eq:estimation} in \eqref{eq:estimation:DA} we obtain \eqref{eq:average:difference:weight}.
\end{proof}
%%%%%%%%%%%%%%%%%%%%%%
%%%%%%%%%%%%%%%%%%%%%%

Another useful estimate is the following lemma.
\begin{lemma}\label{lem:difference:average:2}
Let $\alpha$, $\beta$ and $\delta$ be multi-indices and let $k,l$ be two-dimensional multi-indices. Provided $sh\leq 1$, we have
\begin{equation}\label{lem:difference:average}
\mathbf{D}_{h}^{k}\mathbf{A}_{h}^{l}\left(\partial^{\beta}\left( r\partial^{\alpha}\rho\right) \right)=\partial_{i}^{k_{i}}\partial_{j}^{k_{j}}\partial^{\beta}\left(r\partial^{\alpha}\rho \right)+h^{2}\mathcal{O}_{\lambda}(s^{|\alpha|}).
\end{equation}
Moreover, for $i\neq j\in \{ 1,\ldots,d\}$ we have
\begin{equation}\label{eq0:lem:difference:average:3}
\mathbf{D}_{h}^{k}\mathbf{A}_{h}^{l}\partial^{\delta}\left(r^{2}\left(\partial^{\alpha}\rho\right)\partial^{\beta}\rho\right)=\partial^{k+\delta}\left(r^{2}(\partial^{\alpha}\rho)\partial^{\beta}\rho \right)+h^{2}\mathcal{O}(s^{|\alpha|+|\beta|}),
\end{equation}
and for $\sigma,\sigma'\in[0,1]$ we get
\begin{align}
\mathbf{D}^{k}_{h}\mathbf{A}_{h}^{l}\partial^{\beta}\left(r(x)\partial^{\alpha}\rho(x+\sigma he_{i
} \right))&=\mathcal{O}_{\lambda}(s^{|\alpha|}),\label{eq1:lem:difference:average}\\
\mathbf{D}_{h}^{k}\mathbf{A}_{h}^{l}\partial^{\delta}\left(r^{2}(x)\left( \partial^{\alpha}\rho(x+\sigma he_{i})\right)\partial^{\beta}\rho(x+\sigma' he_{j})\right)&=\mathcal{O}_{\lambda}(s^{|\alpha|+|\beta|}).\label{eq2:lem:difference:average}
\end{align}
\end{lemma}

\begin{proof} The method of proof is analogous to that used in demonstrating Proposition \ref{prop:difference:average}. Specifically, Corollary \ref{cor:difference:average} is applied to $\partial^{\beta}(r\partial^{\alpha}\rho)$, followed by the application of Lemma \ref{lem:derivative:wrt:x} to establish result \eqref{lem:difference:average}. In a similar fashion, equations \eqref{eq0:lem:difference:average:3} and \eqref{eq2:lem:difference:average} are derived using the same approach, with the conclusion utilizing Corollary \ref{Cor:derivative:wrt:x:2r}.
\end{proof}
%%%%%%%%%%%%%%%%%%%%%%%%%%%
%%%%%%%%%%%%%%%%%%%%%%%%%%%
Now, we state the following Theorem that provides us an estimate for several computations of the discrete average and derivative operators applied to our Carleman weight function. 
\begin{theorem}\label{theo:weight:estimates}
Let $\alpha$ be a multi-index and $l,k,n,m,p$ and $q$ be two-dimensional multi-indices. For $i,j\in\{ 1,\ldots, d\}$, and for $sh\leq 1$, we have
\begin{equation}\label{eq1:theo:weight:estimates}
\begin{aligned}
       \mathbf{A}_{h}^{l}\mathbf{D}_{h}^{k}\partial^{\alpha}(r\mathbf{A}_{h}^{m}\mathbf{D}_{h}^{n}\rho)=&\partial^{k+\alpha}(r\partial^{n}\rho)+s^{|\alpha|+n_{i}}\mathcal{O}_{\lambda}((sh)^{2})\\
       &+s^{|\alpha|+n_{j}}\mathcal{O}_{\lambda}((sh)^{2})+s^{|n|}\mathcal{O}_{\lambda}((sh)^{2})
\end{aligned}       
\end{equation}
and 
\begin{equation}\label{eq2:theo:weight:estimates}
\begin{aligned}
\mathbf{A}_{h}^{l}\mathbf{D}_{h}^{k}\partial^{\beta}(r^{2}\mathbf{A}^{p}_{h}\mathbf{D}^{q}_{h}(\partial^{\alpha}\rho)\mathbf{A}_{h}^{m}\mathbf{D}_{h}^{n}\rho)=&\partial^{k+\beta}(r^{2}(\partial^{q+\alpha}\rho)\partial^{n}\rho)\\
&+s^{|\alpha+n+q|}\mathcal{O}_{\lambda}((sh)^{2}).
\end{aligned}
\end{equation}
\end{theorem}
\begin{proof}
From Corollary \ref{cor:difference:average} we write
\begin{equation*}
    r\mathbf{A}_{h}^{m}\mathbf{D}_{h}^{n}\rho= r\partial^{n}(\mathbf{A}_{h}^{m}\rho)+rR_{D^{n_{j}}_{j}}(\mathbf{A}_{h}^{m}\rho)+rR_{D^{n_{i}}_{i}}(\mathbf{A}_{h}^{m}\rho)+rR_{\mathbf{D}_{h}^{n}}(\mathbf{A}_{h}^{m}\rho).
\end{equation*}
Then, applying the operator $\mathbf{A}_{h}^{l}\mathbf{D}_{h}^{k}\partial^{\alpha}$ we have
\begin{align}\label{eq:prop:difference}
    \mathbf{A}_{h}^{l}\mathbf{D}_{h}^{k}\partial^{\alpha}  (r\mathbf{A}_{h}^{m}\mathbf{D}_{h}^{n}\rho)=&  \mathbf{A}_{h}^{l}\mathbf{D}_{h}^{k}\partial^{\alpha}(r\partial^{n}(\mathbf{A}_{h}^{m}\rho))+ \mathbf{A}_{h}^{l}\mathbf{D}_{h}^{k}\partial^{\alpha}(rR_{D^{n_{j}}_{j}}(\mathbf{A}_{h}^{m}\rho))\notag\\
    &+ \mathbf{A}_{h}^{l}\mathbf{D}_{h}^{k}\partial^{\alpha}(rR_{D^{n_{i}}_{i}}(\mathbf{A}_{h}^{m}\rho))+ \mathbf{A}_{h}^{l}\mathbf{D}_{h}^{k}\partial^{\alpha}(rR_{\mathbf{D}_{h}^{n}}(\mathbf{A}_{h}^{m}\rho)).
\end{align}
Let us focus on the term $\mathbf{A}_{h}^{l}\mathbf{D}_{h}^{k}\partial^{\alpha}(r\partial^{n}(\mathbf{A}_{h}^{m}\rho))$. We notice that, for this term, using the first equation from Corollary \ref{cor:difference:average} applied to $\mathbf{A}_{h}^{m}\rho$ gives
\begin{align*}
    \mathbf{A}_{h}^{l}\mathbf{D}_{h}^{k}\partial^{\alpha}(r\partial^{n}(\mathbf{A}_{h}^{m}\rho))=&\mathbf{A}_{h}^{l}\mathbf{D}_{h}^{k}\partial^{\alpha}(r\partial^{n}\rho)+\mathbf{A}_{h}^{l}\mathbf{D}_{h}^{k}\partial^{\alpha}(r\partial^{n}R_{A_{i}^{m_{i}}}(\rho))\\
    &+\mathbf{A}_{h}^{l}\mathbf{D}_{h}^{k}\partial^{\alpha}(r\partial^{n}R_{A_{j}^{m_{j}}}(\rho))+\mathbf{A}_{h}^{l}\mathbf{D}_{h}^{k}\partial^{\alpha}(r\partial^{n}R_{\mathbf{A}_{h}^{m}}(\rho)).
    \end{align*}
    Thus, applying Proposition \ref{prop:weight} to the first term from the right-hand side above, and then \eqref{eq1:lem:difference:average} from Lemma \ref{lem:difference:average:2} we can estimate 
    \begin{align*}
    \mathbf{A}_{h}^{l}\mathbf{D}_{h}^{k}\partial^{\alpha}(r\partial^{n}(\mathbf{A}_{h}^{m}\rho))=&\partial^{k+\alpha}(r\partial^{n}\rho)+s^{|\alpha|+n_{i}}\mathcal{O}_{\lambda}((sh)^{2})+s^{|\alpha|+n_{j}}\mathcal{O}_{\lambda}((sh)^{2})\\
    &+s^{|n|}\mathcal{O}_{\lambda}((sh)^{2}).
    \end{align*}
    Similarly, we can follow the same steps above to estimate the last three terms of \eqref{eq:prop:difference}, which are of order $\mathcal{O}_{\lambda}((sh)^{2})$ and $\mathcal{O}_{\lambda}((sh)^{4})$, and conclude the proof of \eqref{eq1:theo:weight:estimates}. To prove \eqref{eq2:theo:weight:estimates}, we note that thanks to Corollary \ref{cor:difference:average} we write
   \begin{equation*}
    r\mathbf{A}_{h}^{m}\mathbf{D}_{h}^{n}\rho= r\partial^{n}(\mathbf{A}_{h}^{m}\rho)+rR_{D^{n_{j}}_{j}}(\mathbf{A}_{h}^{m}\rho)+rR_{D^{n_{i}}_{i}}(\mathbf{A}_{h}^{m}\rho)+rR_{\mathbf{D}_{h}^{n}}(\mathbf{A}_{h}^{m}\rho)
\end{equation*}
and 
\begin{equation*}
    r\mathbf{A}_{h}^{p}\mathbf{D}_{h}^{q}(\partial^{\alpha}\rho)= r\partial^{q}(\mathbf{A}_{h}^{p}\partial^{\alpha}\rho)+rR_{D^{q_{j}}_{j}}(\mathbf{A}_{h}^{p}\partial^{\alpha}\rho)+rR_{D^{q_{i}}_{i}}(\mathbf{A}_{h}^{p}\partial^{\alpha}\rho)+rR_{\mathbf{D}_{h}^{q}}(\mathbf{A}_{h}^{p}\partial^{\alpha}\rho).
\end{equation*}
Thus, multiplying the above expressions by $\mathbf{A}_{h}^{m}\mathbf{D}_{h}^{n}\rho$ and arguing as in the proof of \eqref{eq1:theo:weight:estimates} the estimate of \eqref{eq2:theo:weight:estimates} follows.
\end{proof}
\begin{remark}
We remark that the estimates in this section remain valid if we interchange $\rho$ and $r$. This is because Lemma \ref{lem:derivative:wrt:x} and Corollary \ref{Cor:derivative:wrt:x:2r} still hold with this change, and we can replace $s$ by $-s$ in these statements. 
\end{remark}

It is worth noting that the previous results can be extended to the case where the weight function is time-dependent. For instance, if we consider a weight function of the form $r(x,t)=e^{\tau\theta(t)\varphi(x)}$, the condition $sh\leq 1$ needs to be replaced by $\tau h(\max_{[0,T]}\theta(t))\leq 1$, which implies $\tau\theta(t)h\leq 1$. This is typically the case for parabolic operators where the Carleman weight function has the form $e^{\tau\theta(t)\varphi(x)}$ and the function $\theta(t)$ is given by 
\begin{equation}\label{theta-delta}
    \theta(t)=\frac{1}{(t+\delta T)(T+\delta T-t)},\quad t\in [0,T],
\end{equation}
for some $0<\delta \leq 1/2$. The parameter $\delta$ is chosen to avoid singularities at time $t=0$ and $t=T$ since $\theta$ must be bounded on the interval $[0,T]$ to obtain an estimate similar to Theorem \ref{theo:weight:estimates} for the time-dependent case. Moreover, we notice that
\begin{equation}\label{eq:theta} 
\max_{t\in [0,T]} \theta(t)=\theta(0)=\theta(T)=\frac{1}{T^{2}\delta(1+\delta)}\leq \frac{1}{T^{2}\delta}, 
\end{equation}
 Thus, the condition $sh\leq 1$ for $\theta$ as in \eqref{theta-delta} is replaced by $sh(T^{2}\delta)^{-1}\leq 1$. This situation is documented in the works \cite{boyer-2014,GC-HS-2021,LMPZ:2023}. Therefore, Theorem \ref{theo:weight:estimates} is a generalization of the estimates presented in those works.
 
\par For the aforementioned Carleman weight function, we obtain, in the case of time derivative, the following asymptotic behavior, which stands for a generalization into an arbitrary dimension of Theorem 2.9 from \cite{CLTP-2022}.  
\begin{theorem}\label{theo:time:derivate}
Let $\alpha$ be a multi-index and let $m,n,j,k$ be two-dimensional multi-indices. Then, provided $\tau h(\delta T^{2})^{-1}\leq 1$, we have
\begin{equation}\label{eq:temporal:estimate}
    \begin{aligned}\partial_{t}\left( r \mathbf{D}_{h}^{n}\mathbf{A}_{h}^{m}\partial^{\alpha}\rho\right)=&\partial_{t}\left(r\partial^{n+\alpha}\rho \right)+Ts^{|\alpha|+n_{i}}\theta(t)\mathcal{O}_{\lambda}((sh)^{2})\\
    &+Ts^{|\alpha|+n_{j}}\theta(t)\mathcal{O}_{\lambda}((sh)^{2})+Ts^{|\alpha|+|n|}\theta(t)\mathcal{O}_{\lambda}((sh)^{2}),
    \end{aligned}
\end{equation}
and
\begin{equation}\label{eq:temporal:estimation}
\begin{split}
    \partial_{t} \mathbf{D}_{h}^{k}\mathbf{A}_{h}^{j}\left(r\mathbf{D}_{h}^{n}\mathbf{A}_{h}^{m}\partial^{\alpha}\rho\right)
    =&T\theta s^{|\alpha|+|n|}\mathcal{O}_{\lambda}(1).
    \end{split}
\end{equation}
\end{theorem}
\begin{proof}
We begin as in the proof of Theorem \ref{theo:weight:estimates}, by using Corollary \ref{cor:difference:average} we have 
\begin{equation*}
\begin{aligned}
   \partial_{t}(r\mathbf{D}_{h}^{n}\mathbf{A}_{h}^{m}\partial^{\alpha}\rho)=&\partial_{t}(r\partial^{n}(\mathbf{A}_{h}^{m}\partial^{\alpha}\rho))+\partial_{t}(rR_{D^{n_{j}}_{j}}(\mathbf{A}_{h}^{m}\partial^{\alpha}\rho))\\
   &+\partial_{t}(rR_{D^{n_{i}}_{i}}(\mathbf{A}_{h}^{m}\partial^{\alpha}\rho))+\partial_{t}(rR_{\mathbf{D}_{h}^{n}}(\mathbf{A}_{h}^{m}\partial^{\alpha}\rho)).
   \end{aligned}
\end{equation*}
Note that applying Corollary \ref{cor:difference:average} to $\partial_{t}(r\partial^{n}(\mathbf{A}_{h}^{m}\partial^{\alpha}\rho))$ we write
\begin{equation}\label{estimate:theo:time1}
\begin{aligned}
    \partial_{t}(r\partial^{n}(\mathbf{A}_{h}^{m}\partial^{\alpha}\rho))=&\partial_{t}(r\partial^{n+\alpha}\rho)+\partial_{t}(rR_{A_{i}^{m_{i}}}(\partial^{n+\alpha}\rho))\\
    &+\partial_{t}(rR_{A_{j}^{m_{j}}}(\partial^{n+\alpha}\rho))+\partial_{t}(rR_{A_{i}^{m_{i}}A_{j}^{m_{j}}}(\partial^{n+\alpha}\rho)).
    \end{aligned}
\end{equation}
Thus, our next task is to estimate the last three terms from the right-hand side above. First, focusing on $\partial_{t}(rR_{A_{i}^{m_{i}}}(\partial^{n+\alpha}\rho))$, thanks to the first estimate from \cite[Proposition 2.14]{boyer-2014} it follows that
\begin{equation}\label{eq:coro:rho}
\partial_{t}(r(x,t)(\partial_{i}^{2}\partial^{n+\alpha}\rho)(x+(m_{i}-2k)e_{i}h\sigma/2,t))=Ts^{|\alpha|+|n|+2}\theta(t)\mathcal{O}_{\lambda}(1).
\end{equation}
Then, with the above estimate and the asymptotic expansion of $R_{A_{i}^{m_{i}}}$ given by Proposition \ref{prop:difference:average} we obtain $\partial_{t}(rR_{A_{i}^{m_{i}}}(\partial^{n+\alpha}\rho))=Ts^{|\alpha|+|n|
}\theta(t)\mathcal{O}_{\lambda}((sh)^{2})$. A similar procedure for the last two terms on the right-hand side of \eqref{estimate:theo:time1} allows to estimate
\begin{align*}
    \partial_{t}(rR_{A_{j}^{m_{j}}}(\partial^{n+\alpha}\rho))=&Ts^{|\alpha|+|n|
}\theta(t)\mathcal{O}_{\lambda}((sh)^{2}),\\
    \partial_{t}(rR_{A_{i}^{n_{i}}A_{j}^{m_{j}}}(\partial^{n+\alpha}\rho))=&Ts^{|\alpha|+|n|
}\theta(t)\mathcal{O}_{\lambda}((sh)^{4}).
\end{align*}
Therefore, the previous last three estimates yield that \eqref{estimate:theo:time1} can be estimated as 
\begin{equation*}
    \partial_{t}(r\partial^{n}(\mathbf{A}_{h}^{m}\partial^{\alpha}\rho))= Ts^{|\alpha|+|n|
}\theta(t)\mathcal{O}_{\lambda}((sh)^{2}).
\end{equation*}

Similarly, we obtain
\begin{align*}
    \partial_{t}(rR_{D^{n_{j}}_{j}}(\mathbf{A}_{h}^{m}\partial^{\alpha}\rho))=&Ts^{|\alpha|+n_{j}}\theta(t)\mathcal{O}_{\lambda}((sh)^{2}),\\
   \partial_{t}(rR_{D^{n_{i}}_{i}}(\mathbf{A}_{h}^{m}\partial^{\alpha}\rho))=&Ts^{|\alpha|+n_{i}}\theta(t)\mathcal{O}_{\lambda}((sh)^{2}),\\
   \partial_{t}(rR_{\mathbf{D}_{h}^{n}}(\mathbf{A}_{h}^{m}\partial^{\alpha}\rho))=&Ts^{|\alpha|+|n|}\theta(t)\mathcal{O}_{\lambda}((sh)^{2}).
\end{align*}
Combining these last four estimates, we establish \eqref{eq:temporal:estimate}.\\
The estimate \eqref{eq:temporal:estimation} can be obtained following a similar strategy.
\end{proof}
\begin{remark}
    Theorem \ref{theo:time:derivate} is a generalization of the second and third estimates from \cite[Proposition 2.14]{boyer-2014} and \cite{HS:2023}.  
\end{remark}

\begin{remark}\label{rmk:1} In the study of inverse problem for  hyperbolic operators \cite{BEO-2015,BE-2013} and \cite{B-B-E-O:2021} 
 is considered the following Carleman weight function 
 \begin{equation}\label{eq:Carleman:wave}
\psi(t,x) = |x - x^{\ast}|^2 - \beta t^2 + c_0, \quad \varphi(t,x) = e^{\mu \psi(t,x)},
\end{equation}
where $x^{\ast}\notin [0,1]^{2}$, and $c_0 > 0$ is such that $\psi \geq 1$ in $[-T,T] \times [0,1]^2$, $\beta\in(0,1)$ and $\mu \geq 1$ is a suitable parameter. Thus, in this setting, we can write
\begin{equation*}
    e^{\lambda\varphi}=e^{\lambda \tilde{\theta}(t) \phi(x)}
\end{equation*}
with $\tilde{\theta}(t):=e^{-\mu\beta t^{2}}$ and $\phi(x):=e^{\mu(|x-x^{\ast}|^{2}+c_{0})}$. Hence, for these weight functions, we can apply Theorem \ref{theo:weight:estimates} to determine its asymptotic behavior since the time variable function is bounded in the interval $[0,T]$. However, this is not the case for Theorem \ref{theo:time:derivate} since the time-dependent factor does not take the same form as $\theta(t)$ in the parabolic case. See also \cite{ZZ:2022} and \cite{ZZ:2023}, where a similar weight function is considered to study the coefficient inverse problems for damped wave and Schr\"{o}dinger equations, respectively. 
\end{remark}
\begin{remark}
In the same direction as the aforementioned remark, in \cite{WZ:2024}, in the study of an inverse problem for a semi-discrete stochastic parabolic operator is considered, a Carleman weight function like \eqref{eq:Carleman:wave}. For this reason, the estimates established above provide a useful tool to extend the results from \cite{WZ:2024} to arbitrary dimensions.
\end{remark}
%%%%%%%%%%%%%%%%%%%%%%%%%%%%%%%
%%%%%%%%%%%%%%%%%%%%%%%%%%%%%%%

Similarly to the previous results, it is possible to obtain an asymptotic expansion considering the time and spatial discrete operators. The main task is to generalize to arbitrary dimensions the asymptotic behavior obtained in \cite{LMPZ:2023}.Now, we will introduce some notations to define the discretization of the time variable. To do this, we recall $\Delta t= T/N$, with $N\in \mathbb{N}$. We define
\begin{align*}
    \mathcal{N}=&\{j\Delta t\,;\, j\in \{ 1,\ldots, N\}\}, \quad \overline{\mathcal{N}}=\mathcal{N}\cup \{0\},\\
    \mathcal{N}^{\ast}=&\{(j-1/2)\Delta t\,;\, j\in \{1,\ldots,N\}\},\quad \overline{\mathcal{N}}^*=\mathcal{N}^{\ast} \cup \{T+\Delta t/2\}. 
\end{align*}
We define the time discrete derivative of a function $y:[0,T]\to \mathbb{R}$ sampled on $\mathcal{N}$ as follows
\begin{align*}
    D_{t} y:=\dfrac{\trp    y-\trm   y}{\Delta t}\quad \text{in } \mathcal{N}^{\ast},
\end{align*}
where the translation in time is defined by $\trpm y(t):=y(t\pm \Delta t/2)$, with $t\in \mathcal{N}$. For the time-discrete operator  we have the following asymptotic behavior
\begin{lemma}\label{lem:time:discrete:weight}
    Provided $\Delta t\tau(T^{3}\delta^{2})^{-1}\leq 1$, we have
    \begin{equation}\label{eq:theta:D:t}
        \trm(\theta^{p}e^{s\varphi})D_{t}(\theta^{-p}e^{-s\varphi})=\trm(\theta^{p}e^{s\varphi}\partial_{t}(\theta^{-p}e^{-s\varphi}))+\mathcal{E}(\Delta t) \mathcal{O}_{\lambda}(1),
    \end{equation}
    where $\mathcal{E}(\Delta t):=\mathcal{E}(\Delta t, p,\delta,\tau,T):=p(p+1)\frac{\Delta t}{\delta^{p+4}T^{2}}+p\frac{\Delta t}{\delta^{p+3}T^{2}}
    +2p\tau\frac{\Delta t}{\delta^{p+4}T^{4}}+\tau^{2} \frac{\Delta t}{\delta^{p+4}T^{6}}+\tau\frac{\Delta t}{\delta^{p+3}T^{4}}$.
\end{lemma}
\begin{proof}
We use the first-order Taylor formula for  $f:=\theta^{-p}e^{-s\varphi}$ to write
\begin{equation}\label{eq:time:taylor}
    \trp D_{t}(f)=\partial_{t}f+\Delta t\int_{0}^{1}(1-\gamma)\partial_{t}^{2}f(t+\gamma \Delta t)d\gamma.
\end{equation}
Let us focus on the integral remainder of the above expression. We have
\begin{equation*}
\begin{aligned}
    \partial_{t}^{2}f=&(p(p+1)\theta^{-p-2}(\partial_{t}\theta)^{2}-p\theta^{-p-1}\partial_{t}^{2}\theta)e^{-s\varphi}+2\tau\varphi d\theta^{-p-1}(\partial_{t}\theta)^{2}e^{-s\varphi}\\
    &+\theta^{-p}\left[ \tau^{2}\varphi^{2}(\partial_{t}\theta)^{2}e^{-s\varphi}-\tau\partial_{t}^{2}\theta\varphi e^{-s\varphi}\right].
      \end{aligned}
\end{equation*}
We notice that there exists a constant $C>0$, only depending on $p$, such that
\begin{equation}\label{ine:theta:d:-d}
    \max_{t\in [0,T]} \theta^{p}(t)\theta^{-p}(t+\gamma\Delta t)\leq \frac{C}{\delta^{p}}.
\end{equation} Indeed, for $p\in\{0,1,2,\ldots\}$ it follows that
\begin{equation}\label{ine:max:theta}
    \max_{t\in[0,T]}\theta^{-p}(t+\gamma\Delta t)\leq   \frac{T^{2p}(1+2\delta)^{2p}}{4^{p}}\leq T^{2p},
    \end{equation}
    and 
    \begin{equation}\label{ine:theta:d}
    \max_{t\in[0,T]} \theta^{p}(t)\leq \frac{C}{\delta^{p}T^{2p}}.
\end{equation} 
So \eqref{ine:theta:d:-d} holds. And we also have
\begin{equation}\label{ine:theta:d:2}
    \max_{t\in[0,T]}\theta^{(j)}(t)\leq \frac{1}{\delta^{j+1}T^{j+2}},\quad j=0,1,\ldots
\end{equation}
Thus, multiplying \eqref{eq:time:taylor} by $\theta^{p}e^{s\varphi}$, and then applying \eqref{ine:theta:d:-d}-\eqref{ine:theta:d:2} to the result we obtain
\begin{equation}
\begin{aligned}
    \theta^{p}e^{s\varphi}\trp D_{t}(\theta^{-p}e^{-s\varphi})=&\theta^{p}e^{s\varphi}\partial_{t}(\theta^{-p}e^{-s\varphi})+p(p+1)\frac{\Delta t}{\delta^{p+4}T^{2}}\mathcal{O}_{\lambda}(1)+p\frac{\Delta t}{\delta^{p+3}T^{2}}\mathcal{O}_{\lambda}(1)\\
    &+2p\tau\frac{\Delta t}{\delta^{p+4}T^{4}}\mathcal{O}_{\lambda}(1)+\tau^{2} \frac{\Delta t}{\delta^{p+4}T^{6}}\mathcal{O}_{\lambda}(1)+\tau\frac{\Delta t}{\delta^{p+3}T^{4}}\mathcal{O}_{\lambda}(1),
\end{aligned}
\end{equation}
where the conditions $\Delta t\tau(\delta^{2} T^{3})^{-1}\leq 1$ and $\|\varphi\|_{\infty}=\mathcal{O}_{\lambda}(1)$ are used to control the exponential term $e^{\tau(\theta(t)-\theta(t+\gamma\Delta t)\varphi)}=\mathcal{O}_{\lambda}(1)$. Then, by shifting the above expression using $\trm$ the estimate \eqref{eq:theta:D:t} holds.
\end{proof}
\begin{remark}
    Following the same strategy, for $p\geq 1$ and $l\geq 0$ be integers and $s(t)=\tau\theta^{p}(t)$, where $\tau\geq 1$. Provided 
    \begin{equation}\label{eq:hypothesis:general}
    \Delta t\,\tau\,(\delta^{p+1}T^{2p+1})^{-1}\leq 1,
    \end{equation}
    we have
    \begin{equation}\label{eq:theta:D:t:general}
        \trm(\theta^{l}e^{s\varphi})D_{t}(\theta^{-l}e^{-s\varphi})=\trm(\theta^{l}e^{s\varphi}\partial_{t}(\theta^{-l}e^{-s\varphi}))+\mathcal{E}(\Delta t) \mathcal{O}_{\lambda}(1),
    \end{equation}
    where $\mathcal{E}(\Delta t):=\mathcal{E}(\Delta t, p, l,\delta,\tau,T)$ is given by
    \begin{equation*}
    \begin{aligned}
    \mathcal{E}(\Delta t):=\Delta t\bigg(&\frac{l(l+1)}{\delta^{l+4}T^{2}}+\frac{l}{\delta^{l+3}T^{2}}+\frac{C(l,p)\,\tau}{\delta^{l+(p-2)^{+}+4}T^{2p+2}}+\frac{p^{2}\tau^{2}}{\delta^{l+2p+2}T^{4p+2}}\bigg),
    \end{aligned}
    \end{equation*}
    with $C(l,p):=|2lp-p(p-1)|+p$ and $(p-2)^{+}:=\max(p-2,0)$.\\
    Moreover, \eqref{eq:theta:D:t:general} generalizes the estimates \cite[Lemma B.4]{BHS:2020}.
\end{remark}

Thus, we will focus on the Carleman weight function of the form $r(x,t):=e^{s(t)\varphi(x)}$ with $\theta$ defined in \eqref{theta-delta}, $\rho=:r^{-1}$ and $s(t)=\tau\theta(t)$.
In \cite{GC-HS-2021} an asymptotic behavior of the Carleman weight function used for the heat equation is established in the one-dimensional setting with limited order in the space-discrete operators.
Moreover, in \cite[Theorem 5]{LMPZ:2023} this limitation is overcome with respect to the order of the spatial operator, although still in the 1-D case. Thus, our task now is to generalize the aforementioned results to arbitrary dimensions. We emphasize that the methodology of the next result adapts the proof of \cite[Lemma A.16]{GC-HS-2021} and \cite[Theorem 5]{LOP-2020}.
\begin{theorem}\label{theo:fully:weight:estimates}
Let $\alpha$ and $m,n$ be multi-index and two-dimensional multi-indices in the space variable, respectively. Provided $\tau h(\delta T^{2})^{-1}\leq 1$, $\Delta t \tau(T^{3}\delta^{2})^{-1}\leq 1/2$ and $\sigma$ bounded, we have
\begin{equation}\label{eq:prop:time:estimate}
    \begin{aligned}
    D_{t}(r\mathbf{D}_{h}^{n}\mathbf{A}_{h}^{m}(\partial^{\alpha}\rho))=&\partial_{t}(r\partial^{n+\alpha}\rho)+T(sh)^{2}s^{|n+\alpha|}\theta\mathcal{O}_{\lambda}(1)\\
    &+\Delta t \trp((sh)^{2}T^{2}s^{|n+\alpha|}\theta^{2})\mathcal{O}_{\lambda}(1)\\
    &+\Delta t T^{2}\trp(s^{|n+\alpha|}\theta^{2})\mathcal{O}_{\lambda}(1).
\end{aligned}
\end{equation}
\end{theorem}
\begin{proof}
   To prove \eqref{eq:prop:time:estimate} we follow a similar strategy developed in the proof of Theorem \ref{theo:weight:estimates}. Thanks to Corollary \ref{cor:difference:average} we write
\begin{equation}\label{estimate:theo:fully}
\begin{aligned}
    r\mathbf{D}_{h}^{n}\mathbf{A}_{h}^{m}\partial^{\alpha}\rho=& r\partial^{n}(\mathbf{A}_{h}^{m}\partial^{\alpha}\rho)+rR_{D^{n_{j}}_{j}}(\mathbf{A}_{h}^{m}\partial^{\alpha}\rho)+rR_{D^{n_{i
}}_{i}}(\mathbf{A}_{h}^{m}\partial^{\alpha}\rho)\\
&+rR_{\mathbf{D}_{h}^{n}}(\mathbf{A}_{h}^{m}\partial^{\alpha}\rho).
    \end{aligned}
\end{equation}
Notice that the first term from the above right-hand side can be written as
\begin{align*}
r\partial^{n}(\mathbf{A}_{h}^{m}\partial^{\alpha}\rho)=&r\partial^{n+\alpha}\rho+rR_{A_{j}^{m_{j}}}(\partial^{\alpha+n}\rho)+rR_{A_{i}^{m_{i}}}(\partial^{\alpha+n}\rho)+rR_{\mathbf{A}_{h}^{m}}(\partial^{\alpha+n}\rho)
\end{align*}
where we have used Corollary \ref{cor:difference:average}.
Then, the first-order Taylor formula in the time variable yields
\begin{equation}\label{eq:almost:estimate}
\begin{aligned}
    D_{t}(r\partial^{n}(\mathbf{A}_{h}^{m}\partial^{\alpha}\rho))=&\partial_{t}(r\partial^{n+\alpha}\rho)+\partial_{t}(rR_{A_{j}^{m_{j}}}(\partial^{n+\alpha}\rho))\\
    &+\partial_{t}(rR_{A_{i}^{m_{i}}}(\partial^{\alpha+n}\rho))+\partial_{t}(rR_{\mathbf{A}_{h}^{m}}(\partial^{\alpha+n}\rho))\\
    &+\Delta t\int_{0}^{1}(1-\gamma)\partial_{t}^{2}\left( r(x,t+\gamma\Delta t)\partial^{n}\mathbf{A}_{h}^{m}(\partial^{\alpha}\rho(x,t+\gamma\Delta t))\right)d\gamma.
    \end{aligned}
\end{equation}
Theorem \ref{theo:time:derivate} enables to estimate
\begin{equation}\label{eq:estimate:R}
\begin{aligned}
\partial_{t}(rR_{A_{j}^{m_{j}}}(\partial^{n+\alpha}\rho))=&T(sh)^{2}s^{|\alpha|+|n|}\theta\mathcal{O}_{\lambda}(1),\\ \partial_{t}(rR_{A_{i}^{m_{i}}}(\partial^{\alpha}\rho))=&T(sh)^{2}s^{|\alpha|+|n|}\theta\mathcal{O}_{\lambda}(1),
\end{aligned}
\end{equation}
and 
\begin{equation}\label{eq:estimate:R:01}
\partial_{t}(rR_{\mathbf{A}_{h}^{m}}(\partial^{\alpha+n}\rho))=T(sh)^{4}s^{|\alpha|+|n|}\theta\mathcal{O}_{\lambda}(1).
\end{equation}
According to the first estimate from \cite[Lemma A.16]{GC-HS-2021} and Proposition \ref{prop:weight}, the last term from the right-hand side of \eqref{eq:almost:estimate} can be estimated as
\begin{equation}\label{integral:remainder}
\begin{split}
    \int_{0}^{1}(1-\gamma)\partial_{t}^{2}\left( r(x,t+\gamma\Delta t)\mathbf{A}_{h}^{m}(\partial^{\alpha}\rho(x,t+\gamma\Delta t))\right)d\gamma=& T^{2}\trp(s^{|n|+|\alpha|}\theta^{2})\mathcal{O}_{\lambda}(1)\\
    &+ \trp((sh)^{2}T^{2}s^{|n|+|\alpha|}\theta^{2})\mathcal{O}_{\lambda}(1).
    \end{split}
\end{equation}
Hence, combining \eqref{eq:estimate:R}-\eqref{integral:remainder}, for \eqref{eq:almost:estimate} we obtain 
\begin{equation}
\begin{aligned}
    D_{t}(r\partial^{n}(\mathbf{A}_{h}^{m}\partial^{\alpha}\rho))=&\partial_{t}(r\partial^{n+\alpha}\rho)+T(sh)^{2}s^{|n|+|\alpha|}\theta\mathcal{O}_{\lambda}(1)\\
    &+\Delta t \trp((sh)^{2}T^{2}s^{|n|+|\alpha|}\theta^{2})\mathcal{O}_{\lambda}(1)\\
    &+\Delta t T^{2}\trp(s^{|n|+|\alpha|}\theta^{2})\mathcal{O}_{\lambda}(1).
    \end{aligned}
\end{equation}

Mimicking the same steps as before, it follows the estimates for the time-discrete operaror applied on the remainder terms of \eqref{estimate:theo:fully}; which yields \eqref{eq:prop:time:estimate}. 
\end{proof}
%%%%%

\section{Carleman estimate}\label{sec:Carleman}

%%%%%%%%%%%%%%%%%%%%

\par The assumption on the Carleman weight function is standard for this problem, since have been used in \cite{boyer-2014,LPP:2026,LLRP:2025,BHLR:2010,LPP:2026b}.\\ \textbf{Assumption:} Let $\overline{\omega_0}\subset \omega$ be an arbitrary fixed subdomain of $\Omega$. Let $\tilde{\Omega}$ be a smooth, open, and connected neighborhood of $\overline{\Omega}$ in $\mathbb{R}^d$. The function $x \mapsto \psi(x)$ belongs to $\mathcal{C}^p(\tilde{\Omega}, \mathbb{R})$, for sufficiently large $p$, and satisfies, for some $c > 0$,
\begin{equation}\label{assumtion:psi}
\psi > 0 \quad \text{in } \widehat{\Omega}, \quad |\nabla \psi| \geq c \quad \text{in } \tilde{\Omega} \setminus \omega_0, \quad \text{and} \quad \partial_{n_i} \psi(x) \leq -c < 0 \quad \text{for } x \in V_{\partial_i \Omega},
\end{equation}
where $V_{\partial_i \Omega}$ is a sufficiently small neighborhood of $\partial_i \Omega$ in $\tilde{\Omega}$, where the outward unit normal $n_i$ to $\Omega$  extends from $\partial_i \Omega$. 
For $\lambda\geq 1$ and $K>\|\psi\|_{\infty}$, we introduce the functions
\begin{align}\label{funcion-peso-2}
\varphi(x)=e^{\lambda\psi(x)}-e^{\lambda K}<0,
\end{align}
 and, for $0<\delta \leq 1/2$,
\begin{equation}
    \theta(t)=\frac{1}{(t+\delta T)(T+\delta T-t)},\quad t\in [0,T].
\end{equation}
Given $\tau\geq 1$, we set
\begin{equation}\label{eq1}
s(t)=\tau\theta(t).
\end{equation}
Applying the asymptotic behavior of the Carleman weight functions given by Theorems \ref{theo:weight:estimates} and \ref{theo:fully:weight:estimates} and following the strategies used in \cite{boyer-2014,BHS:2020,GC-HS-2021,LMPZ:2023,LLRP:2025} we prove a fully-discrete Carleman estimate in arbitrary dimensions. Indeed, denoting  $\mathcal{P}q:=D_{t}q+\sum_{i=1}^{d}D_{i}(\gamma_{i}D_{i}q)$ and 
\begin{equation*}
\begin{aligned}
\mathcal{J}(q):=&\int_{\mathcal{M}\times\mathcal{N}}\trm(s^{-1}r^{2})|D_{t}q|^{2}+\sum_{i=1}^{d}\int_{\mathcal{M}\times\mathcal{N}^{\ast}}s^{-1}r^{2}|D_{i}(\gamma_{i}D_{i}q)|^{2}+\sum_{i=1}^{d}\int_{\mathcal{M}_{i}^{\ast}\times\mathcal{N}^{\ast}}s\gamma_{i}\phi\,r^{2}|D_{i}q|^{2}\\
&+\int_{\mathcal{M}\times\mathcal{N}^{\ast}}s^{3}r^{2}\phi^{3}|\nabla_{\Gamma}\psi|^{4}|q|^{2}+\sum_{i=1}^{d}\int_{\mathcal{M}_{i}^{\ast}\times\mathcal{N}^{\ast}}s\gamma_{i}\phi\,r^{2}|A_{i}D_{i}q|^{2}
\end{aligned}
\end{equation*}
we have
\begin{theorem}\label{thm:Carleman:q}
Let $\omega$ be an open subset of $\mathcal{M}$, $\omega_{0}\Subset \omega$ a nonempty open set, $\psi$ satisfying $|\nabla\psi|\geq c>0$ in $\Omega\setminus \overline{\omega}_{0}$, and $\varphi=e^{\lambda\psi}-e^{\lambda K}<0$. For $\lambda\geq 1$ sufficiently large, there exist $C>0$, $\tau_{0}\geq 1$, $\epsilon_{0}>0$ depending on $\omega$, $\omega_{0}$, $T$, $\lambda$, $\Gamma$ such that
\begin{equation}\label{eq:Carleman:q}
\begin{aligned}
&\mathcal{J}(q)+\sum_{i=1}^{d}\int_{\partial_{i}\mathcal{M}\times\mathcal{N}^{\ast}}s\gamma_{i}\phi r^{2}\,t_{r}^{i}(|D_{i}q|^{2})+\sum_{i=1}^{d}\int_{\partial_{i}\mathcal{M}\times\mathcal{N}^{\ast}}s^{3}r^{2}\phi^{3}t_{r}^{i}(A_{i}|q|^{2})+\int_{\partial\mathcal{M}\times\mathcal{N}}(\trm r)^{2}|D_{t}q|^{2}\\
&\leq C_{\lambda}\bigg(\int_{\mathcal{M}\times\mathcal{N}}(\trm r)^{2}|\mathcal{P}q|^{2}+\int_{\omega\times\mathcal{N}^{\ast}}s^{3}r^{2}\phi^{3}|q|^{2}+h^{-2}\int_{\partial\mathcal{M}\times\partial\mathcal{N}}|\trp(rq)|^{2}+h^{-2}\int_{\mathcal{M}\times\partial\mathcal{N}}|\trp(rq)|^{2}\bigg),
\end{aligned}
\end{equation}
for all $\tau\geq\tau_{0}(T+T^{2})$, $0<h<h_{0}$, $\Delta t>0$ and $0<\delta\leq 1/2$ such that
\begin{equation}\label{eq:CFL:conditions}
\frac{\tau h}{\delta T^{2}}\leq\epsilon_{0}\quad\text{and}\quad\frac{\tau^{4}\Delta t}{\delta^{4}T^{6}}\leq\epsilon_{0}.
\end{equation}
\end{theorem}
This Carleman estimate implies an relaxed observability inequality for the system \eqref{eq:adjoint}.
\begin{theorem}\label{thm:observability}
Let $T>0$, $\mu\geq 1$, and $h_{0}>0$ given by Theorem~\ref{thm:Carleman:q}. There exist constants $C_{0},C,C_{1},C_{\mathrm{obs}}>0$ such that for all $ q_{T}\in L^{2}_{h}(\mathcal{M})$, the solution $q\in C(\overline{\mathcal{M}}\times\overline{\mathcal{N}}^{\ast})$ of \eqref{eq:adjoint} satisfies
\begin{equation}\label{eq:observability}
\|\trp q(0)\|^{2}_{L^{2}_{h}(\mathcal{M})}\leq C_{\mathrm{obs}}\bigg(\int_{\omega\times\mathcal{N}^{\ast}}|q|^{2}+e^{-\frac{C_{1}}{h^{\min\{\mu/4,1\}}}}\|q_{T}\|^{2}_{L^{2}_{h}(\mathcal{M})}\bigg),
\end{equation}
provided $h\leq\min\{h_{0},h_{1}\}$ and $\Delta t\leq\min\{T^{-2}h^{\mu},(4\|a\|_{L^{\infty}_{h}})^{-1}\}$, where
\begin{equation*}
h_{1}:=C_{0}\bigg(1+\frac{1}{T}+\|a\|_{L^{\infty}_{h}}^{2/3}\bigg)^{-\max\{1,4/\mu\}}\quad\text{and}\quad C_{\mathrm{obs}}:=e^{C\big(1+\frac{1}{T}+\|a\|_{L^{\infty}_{h}}^{2/3}+T\|a\|_{L^{\infty}_{h}}\big)}.
\end{equation*}
\end{theorem}
%%%%%%%%%%%%%%%%%%%%
Its proof mimics the steps developed in \cite[section 4]{LMPZ:2023}. Indeed, let $q$ be a solution of \eqref{eq:adjoint}. Then, for all $t\in\mathcal{N}$,
\begin{equation}\label{eq:energy:monotonicity}
\|\trp q(0)\|^{2}_{L^{2}_{h}(\mathcal{M})}\leq\|\trp q(t)\|^{2}_{L^{2}_{h}(\mathcal{M})},
\end{equation}
provided $\Delta t\,\|a\|_{L_{h}^{\infty}}<1/2$. Inequality follows applying the discrete Gronwall inequality (Lemma~2 from \cite{LMPZ:2023}) with $\gamma=2\|a\|_{L_{h}^{\infty}}$ and provided $2\|a\|_{L_{h}^{\infty}}\,\Delta t<1$. Moreover, defining $\eta(t)=1$ for $t\in\mathcal{N}^{\ast}\cap(T/4,3T/4)$ and $\eta(t)=0$ otherwise, applying the Carleman estimate \eqref{eq:Carleman:q} and using \eqref{eq:energy:monotonicity} yield 
\begin{equation}\label{eq:initial:from:energy}
\|\trp q(0)\|^{2}_{L^{2}_{h}(\mathcal{M})}\leq\frac{C}{T}\int_{\mathcal{N}^{\ast}}\eta(t)\|\trp q(t)\|^{2}_{L^{2}_{h}(\mathcal{M})}\leq\frac{C}{T}\int_{\mathcal{M}\times\mathcal{N}^{\ast}}\eta(t)|q|^{2}.
\end{equation}
Implying
\begin{equation}\label{eq:LHS:controls:initial}
\tau^{3}\int_{\mathcal{M}\times\mathcal{N}^{\ast}}\theta^{3}e^{2\tau\theta\varphi}\phi^{3}|q|^{2}\geq C_{T}\,e^{-\frac{C\tau}{T^{2}}}\|\trp q(0)\|^{2}_{L^{2}_{h}(\mathcal{M})}.
\end{equation}
For the boundary terms one can prove
\begin{equation}\label{eq:time:boundary:bound}
h^{-2}\int_{\mathcal{M}\times\partial\mathcal{N}}|\trp(rq)|^{2}\leq C\,h^{-2}\,e^{-\frac{4\tau k_{0}}{\delta T^{2}}+C\|a\|_{L_{h}^{\infty}}T}\|\trp q(T)\|^{2}_{L^{2}_{h}(\mathcal{M})}.
\end{equation}
Since $\theta(0)=\theta(T)=\frac{1}{\delta T(T+\delta T)}\leq\frac{2}{\delta T^{2}}$ (for $\delta\leq 1/2$), and $\varphi<0$ with $k_{0}:=\min(-\varphi)>0$:
\begin{equation*}
e^{2\tau\theta(0)\varphi}\leq e^{-\frac{4\tau k_{0}}{\delta T^{2}}},\qquad e^{2\tau\theta(T)\varphi}\leq e^{-\frac{4\tau k_{0}}{\delta T^{2}}}.
\end{equation*}
%%%%%%%%%%%%%%%%%%
\section{Concluding remarks and perspectives}\label{sec:concluding}
In this work, we established the asymptotic behavior of the Carleman weight functions in arbitrary dimensions for spatial and time discrete operators, which are fundamental tools in developing Carleman estimates for semi-discrete and fully discrete operators. Recent works, such as \cite{GC-HS-2021} and \cite{LMPZ:2023}, have used Carleman estimates to analyze the controllability of fully discrete 1-D heat equations with Dirichlet and dynamic boundary conditions respectively. Hence, this work extends the results to arbitrary dimensions presented in \cite{GC-HS-2021}. For this reason it seems promising using the result presented in this paper as Theorems \ref{theo:weight:estimates} and \ref{theo:fully:weight:estimates} to extend the results obtained in \cite{LMPZ:2023}. Moreover, as is mentioned in remark \ref{rmk:1}, the asymptotic behavior also holds for the weight functions utilized in Carleman estimate for semi-discrete hyperbolic operators then a similar approach as the inverse problems could be adapted into arbitrary dimensions for these semi-discrete operators.
\par In a similar direction, another unexplored research topic where Carleman estimates are fundamental is the discrete or semi-discrete inverse problem for parabolic operators. In recent years, results for semi-discrete hyperbolic operators \cite{BEO-2015,BE-2013} have provided stability estimates for inverse problems. Nevertheless, there is only one stability estimate results for semi-discrete deterministic parabolic operators \cite{LLRP:2025}. Further, \cite{Wu_2024} considers an inverse problem for a semi-discrete stochastic parabolic equation in the one-dimensional setting. Thus, the results presented in this work are important tools for studying those semi-discrete inverse problems, as Carleman estimates are the main tool to achieving these results, see remark \ref{rmk:1}. Indeed, in \cite{LPP:2026b} are obtained to three Carleman estimates for semi-discrete stochastic parabolic operator in arbitrary dimension by using Theorems \ref{theo:weight:estimates} and \ref{theo:fully:weight:estimates}. These Carleman estimates are applied to study inverse random source and Cauchy problems, see \cite[Theorem 1.1]{LPP:2026b} and \cite[Theorem 1.3]{LPP:2026b}.
\par Recently, several works have been considered the controllability problems for semi-discrete stochastic partial differential equations. In \cite{zhao:2024}, \cite{WZ:2025} and \cite{WZ:2024} is established the $\phi$-null controllability of second and fourth-order semi-discrete stochastic parabolic operators, (see \cite{boyer-canum} for details of $\phi$-null controllability for semi-discrete systems). Moreover, in \cite{LL:2026} the authors obtain $\phi$-null controllability resutls for a cascade system of backward stochastic semi-discrete fourth- and second-order parabolic. Concerning the controllability issue, the task is to establish a relaxed observability inequality through a  Carleman estimate for the respective semi-discrete backward adjoint system, then to obtain $\phi$-null controllability result for its respective forward semi-discrete stochastic system. Using this methodology, it is obtained in \cite{zhao:2024}, \cite{WZ:2025} and  \cite{WZ:2024} controllability results for a spatial semi-discrete approximation of a 1-D stochastic parabolic operators. It is worth noting that spatial terms are key to obtaining the Carleman estimates present in those works. For this reason, the estimates related to the Carleman weight function proved in the present paper are fundamental in extending the aforementioned results to arbitrary dimensions. Moreover, due to the difficulty of managing the estimation on the Carleman weight function, there are only two works concerning a higher order operator in 1-D, see \cite{CLTP-2022} and \cite{WZ:2024} for controllability results of a deterministic and stochastic fourth-order operator, respectively. This difficulties, by using Theorems \ref{theo:weight:estimates}, are overcome in \cite{LPP:2026} where the authors obtain $\phi$-controllability result for semi-discrete stochastic parabolic operators in arbitrary dimension which generalize the results obtained in \cite{zhao:2024}.
\par In turn, in \cite{ZY-2024} is obtained a stability estimate of the discrete Calderon's problem with partial data measured on an arbitrary portion of the boundary where a suitable Carleman estimate is the main tool. This is an improvement of the stability result presented in \cite{LDOP-2021} and also needs to estimate repeated application of the discrete spatial operator on the Carleman weight function to obtain a weaker version, with respect to its continuous framework, of the unique continuation property. In the same direction, in \cite{FB:R:R-2021} is obtained a discrete version of the three spheres inequality, and in \cite{FB:R:R:2024} is studied the unique continuation properties of the discrete fractional Laplacian operator both of them via discrete Carleman estimate where their approach mimics the methodology from \cite{ervedoza-2011}. Hence several computations related with trigonometric identities are needed in the development of the discrete Carleman estimate in the works \cite{FB:R:R-2021,ervedoza-2011,LDOP-2021}. Thus, the estimates from Proposition \ref{prop:weight} and Theorem \ref{theo:weight:estimates} give an alternative way to develop of Carleman estimates for discrete elliptic operators. 

\section*{Acknowledgment}
This work was partially supported by the Vicerrector\'ia de Investigaci\'on y Postgrado at Universidad del B\'io-B\'io, project IN2450902; and FONDECYT Grant 11250805. Part of this work was completed during the author's research stay at Escuela de Matem\'aticas, Universidad Industrial de Santander, Colombia. The author extends heartfelt thanks to Professor Juan L\'opez-R\'ios for his kind hospitality.
%%%%%%%%%%%%%%%%%%%%%%
%%%%%%%%%%%%%%%%%%%
%%%%%%%%%%%%%%%%%%%


\section{Proof of the Carleman estimate}\label{sec:proof:carleman}
%%%%%%%%%%%%%%%%%%%%%
%%%%%%%%%%%%%%%%%%%%% 
%-----------
\subsection{Conjugated operator}
%------------
For simplicity of notation we write $r:=e^{s\varphi}$ and $\rho:=r^{-1}$, where our weight function is defined as $\varphi=e^{\lambda\psi}-e^{\lambda K}$, for $s\geq 1$, with $\psi\in C^{k}$ for $k$ sufficiently large and $\lambda\geq 1$. The proof of some technical results can be found in the section 5. 
\par By regarding this weight function, the first step is to consider the change of variable $q=\rho z$. Our first task is to split the conjugate operator $\trm r\mathcal{P}(\rho z):=-\trm rD_{t}(\rho z)-\trm r\sum_{i=1}^{d}\trm D_{i}\left(\gamma_{i}D_{i}(\rho z)\right)$ into simple terms that we will estimate separately. By \eqref{eq:discreteleibniz}, we have
 \begin{equation}\label{eq:conjugate}
\begin{split}
\trm r\sum_{i=1}^{d}\trm D_{i}\left(\gamma_{i}D_{i}(\rho z)\right)=&A_{1}z+A_{2}z+B_{1}z+C_{h}z,
\end{split}
\end{equation}
where $\displaystyle A_{1}z:=\sum_{i=1}^{d}\trm(rA_{i}^{2}\rho\,D_{i}(\gamma_{i}D_{i}z))$, $\displaystyle B_{1}z:=2\sum_{i=1}^{d}\gamma_{i}\trm(rD_{i}A_{i}\rho D_{i}A_{i}z)$,\\ $\displaystyle A_{2}z:=\sum_{i=1}^{d}\gamma_{i}\trm(rD_{i}^{2}\rho A_{i}^{2}z),$
and
\begin{equation*}
    \begin{split}
        C_{h}z:=&\sum_{i=1}^{d}h_{i}\mathcal{O}(1)\trm(rD_{i}^{2}\rho A_{i}^{2}z)+\sum_{i=1}^{d}h_{i}\mathcal{O}(1)\trm(rD_{i}A_{i}\rho D_{i}A_{i}z)+\sum_{i=1}^{d}D_{i}\gamma_{i}\trm(A_{i}D_{i}\rho A_{i}^{2}z)\\
        &+\sum_{i=1}^{d}\frac{h^{2}_{i}}{4}D_{i}\gamma_{i}\trm(D_{i}^{2}\rho D_{i}A_{i}z)+\sum_{i=1}^{d}\frac{h^{2}_{i}}{4}\trm(D_{i}A_{i}\rho D_{i}^{2}z).
    \end{split}
\end{equation*}
Now, using that $\trp z=\trm z+\Delta t D_{t}z$ and Lemma \ref{lem:time:discrete:weight} with $p=0$, provided $\Delta t \tau (T^3 \delta^2)^{-1}\leq 1$ it follows that
\begin{equation}\label{eq:time:conjugate}
\begin{split}
    \trm(r)\,D_{t}(\rho z)
    =&D_{t}z-\tau\varphi\trm(\theta'z)+C_{\Delta t}z,
    \end{split}
\end{equation}
with $C_{\Delta t}z:=\Delta t\left(\frac{\tau}{\delta^{3}T^{4}}+\frac{\tau^{2}}{\delta^{4}T^{6}} \mathcal{O}_{\lambda}(1)\right)\trp z+\tau\Delta t\trm(\theta')\varphi D_{t}z$. Now, we set $A_{3}z:=-\tau\varphi\trm(\theta'z)$, $B_{2}z:=-2\trm(s\Delta_{\Gamma}\phi z)$ and $B_{3}z:=D_{t}z$, where $\Delta_{\Gamma}\phi=\sum_{i=1}^{d}\gamma_{i}\partial_{x_{i}}^{2}\phi$. Thus, combining \eqref{eq:conjugate} and \eqref{eq:time:conjugate} we have the following identity in $\mathcal{M}\times\mathcal{N}$
\begin{equation}\label{eq:innner:product}
    Az+Bz=Rz,
\end{equation}
where $Az=A_{1}z+A_{2}z+A_{3}z$, $Bz:=B_{1}z+B_{2}z+B_{3}z$ and $Rz:=-\trm r\mathcal{P}(\rho z)-C_{h}z-C_{\Delta t}z+B_{2}z$. Hence, we take $L_{h}^{2}(\mathcal{M}\times\mathcal{N})$ in \eqref{eq:innner:product} yields 
\begin{equation}\label{eq:inner:product:2}
    \left\| Az\right\|^{2}+\left\| Bz\right\|^{2}+2\int_{\mathcal{M}\times\mathcal{N}}Az\cdot Bz=\left\| Rz\right\|^{2}.
\end{equation}
 %%%%%%%%%%%%%%%%%%%%%%%%%%%%%
%%%%%%%%%%%%%%%%%%%%%%%%%%%%
The main task to obtain a Carleman estimate is to estimate the following term 
\begin{equation}\label{conmu}
    \int_{\mathcal{M}\times\mathcal{N}}Az\cdot B z=\sum_{i,j=1}^{3}\int_{\mathcal{M}\times\mathcal{N}}A_{i}z\cdot B_{j}z=:\sum_{i,j=1}^{3}I_{ij}. 
\end{equation}
 For every term of (\ref{conmu}) we present its respective estimate in section \ref{sec:cross}. Collecting all these estimation, we get
 \begin{equation}\label{eq:combined:recall}
\begin{aligned}
2\int_{\mathcal{M}\times\mathcal{N}}Az\cdot Bz\geq\;&s^{3}\lambda^{4}\int_{\mathcal{M}\times\mathcal{N}^{\ast}}\phi^{3}|\nabla_{\Gamma}\psi|^{4}\,|z|^{2}\\
&+s\lambda^{2}\sum_{i=1}^{d}\int_{\mathcal{M}_{i}^{\ast}\times\mathcal{N}^{\ast}}\gamma_{i}\phi\Big[2\gamma_{i}(\partial_{i}\psi)^{2}+|\nabla_{\Gamma}\psi|^{2}\Big]|D_{i}z|^{2}\\
&-2s\lambda\sum_{i=1}^{d}\int_{\mathcal{M}_{i}^{\ast}\times\mathcal{N}^{\ast}}\gamma_{i}\phi\,\partial_{i}\psi\,\partial_{i}\gamma_{i}\,|D_{i}z|^{2}\\
&+C_{\lambda}\Delta t\sum_{i=1}^{d}\int_{\mathcal{M}_{i}^{\ast}\times\mathcal{N}}\gamma_{i}\,|D_{it}^{2}z|^{2}+\sum_{i\ne j}V_{11}^{ij}-\mathcal{X}-\mathcal{Y},
\end{aligned}
\end{equation}
where $\mathcal{X}$ and $\mathcal{Y}$ are sub-leading terms that can be absorbed increasing the parameters $\lambda$ and $\tau$ if is necessary. At this point the strategy to follows is from \cite[Lemma 2.12]{GC-HS-2021}. To remove the local term $|D_{i}z|^{2}$ from the left-hand side it follows from \cite[Lemma 2.16]{GC-HS-2021}. Moreover, returning to the original variable mimics \cite[Lemma 9]{LMPZ:2023}, and to incorporate the second-order spatial operator we apply the steps used in \cite[Theorem 1.2]{LLRP:2025}.
%%%%%%%%%%%%%%%%%%%%%%%%%%%%%%%%%%%%%%%%
\section{Estimate for the cross-product}\label{sec:cross}
In this section we provided all the estimate to obtain a bound for the cross-product $\int_{\mathcal{M}\times\mathcal{N}}Az\cdot Bz$. We begin recalling the  definition of the respective $I_{ij}$ for $i,j=1,2,3$, then usually a time-shifting allows to integrate on $\mathcal{M}\times\mathcal{N}^{\ast}$. In higher dimension this is the hardest part in the development of a semi(fully)-discrete Carleman estimate, since in every estimate of $I_{ij}$ is needed to obtain an asymptotic behavior of the Carleman weight function, then several additional estimation are needed to obtain a bound for $I_{ij}$. For this reason, it becomes easier by using Theorem \ref{theo:weight:estimates} and \ref{theo:fully:weight:estimates} since they contain all the possible configuration to be estimated in the proof of a Carleman estimate. Thus, these Theorems not only reduce the computation also enable one to obtain Carleman estimate in higher-dimensions and higher-order. At the end of each section we give a remark about the role of the respective estimate in the Carleman estimate .
\subsection{Estimate of $I_{11}$ }
%%%%%%%%%%%%%%%%%%%%%%%%%%%%%%%%%%%%%%%%
We set $\beta_{11}^{ij}:=\gamma_{j}r^{2}A_{i}^{2}\rho D_{j}A_{j}\rho$, then recalling that
\begin{equation*}
    \begin{split}
        A_{1}z&:=\sum_{i=1}^{d}\trm(rA_{i}^{2}\rho\,D_{i}(\gamma_{i}D_{i}z))\, \text{ and }\,B_{1}z:=2\sum_{i=1}^{d}\gamma_{i}\trm(rD_{i}A_{i}\rho D_{i}A_{i}z),
    \end{split}
\end{equation*}
we have after a shift with respect to the discrete-time variable
\begin{equation*}
    \begin{split}
     I_{11}=\sum_{i,j= 1}^{d}2\int_{\mathcal{M}\times\mathcal{N}} \trm(\beta_{11}^{ij}D_{i}(\gamma_{i}D_{i}z)\,D_{j}A_{j}z)=&\sum_{i,j=1}^{d}2\int_{\mathcal{M}\times\mathcal{N}^{\ast}} \beta_{11}^{ij}D_{i}(\gamma_{i}D_{i}z)\,D_{j}A_{j}z\\
    =:&\sum_{i,j=1}^{d}I_{11}^{ij}.
    \end{split}
\end{equation*}
\par Let us first consider the case $i=j$. Using the product rule \eqref{D:product} we have
\begin{equation*}
\begin{split}
    I_{11}^{ii}=&2\int_{\mathcal{M}\times\mathcal{N}^{\ast}} \beta_{11}^{ii}D_{i}\gamma_{i}|D_{i}A_{i}z|^{2}+\beta_{11}^{ii}A_{i}\gamma_{i}D_{i}^{2}zD_{i}A_{i}z\\
    =&2\int_{\mathcal{M}\times\mathcal{N}^{\ast}} \beta_{11}^{ii}D_{i}\gamma_{i}|D_{i}A_{i}z|^{2}+\int_{\mathcal{M}\times\mathcal{N}^{\ast}} \beta_{11}^{ii}A_{i}\gamma_{i}D_{i}(|D_{i}z|^{2})\\
    =:&I_{11}^{ii(a)}+I_{11}^{ii(b)}.
    \end{split}
\end{equation*}
Applying a discrete integration by parts with respect to the difference operator $D_{i}$ given by \cite[Lemma 2.2, equation (15)]{LDOP-2021} on the integral $I_{11}^{ii(b)}$, it follows that
\begin{equation*}
    I_{11}^{ii(b)}=-\int_{\mathcal{M}^{\ast}_{i}\times\mathcal{N}^{\ast}} D_{i}(\beta_{11}^{ii}A_{i}\gamma_{i})|D_{i}z|^{2}+\int_{\partial_{i}\mathcal{M}\times\mathcal{N}^{\ast}}\beta_{11}^{ii}A_{i}\gamma_{i}t_{r}^{i}(|D_{i}z|^{2})n_{i}.
\end{equation*}
Now, using Theorem \ref{theo:weight:estimates} we estimate $\beta_{11}^{ii}=-    s\lambda\gamma_{i}\phi\partial_{i}\psi+s\mathcal{O}_{\lambda}((sh)^{2})$ and $D_{i}(\beta_{11}^{ii}A_{i}\gamma_{i})=-s\lambda^{2}\gamma_{i}^{2}\phi(\partial_{i}\psi)^{2}+s\lambda\mathcal{O}_{\lambda}(1)+s\mathcal{O}_{\lambda}((sh)^{2})$. Then for $I_{11}^{ii(b)}$ we obtain
\begin{equation}\label{eq:J11b}
\begin{aligned}
I_{11}^{ii(b)}=&\int_{\mathcal{M}^{\ast}_{i}\times\mathcal{N}^{\ast}}\left(s\lambda^{2}\gamma_{i}^{2}\phi(\partial_{i}\psi)^{2}+s\lambda\mathcal{O}_{\lambda}(1)+s\mathcal{O}_{\lambda}((sh)^{2})\right)|D_{i}z|^{2}\\
&+\int_{\partial_{i}\mathcal{M}\times\mathcal{N}^{\ast}}(-    s\lambda\gamma_{i}\phi\partial_{i}\psi+s\mathcal{O}_{\lambda}((sh)^{2}))t_{r}^{i}(|D_{i}z|^{2})n_{i}.
\end{aligned}
\end{equation}
In turn, using the identity $A_{i}(|D_{i}z|^{2})=|D_{i}A_{i}z|^{2}+\frac{h^{2}}{4}|D_{i} ^{2}z|^{2}$ for $i\in\{1,\ldots,d\}$, we can rewrite the integral $I_{11}^{ii(a)}$ as
\begin{equation*}
    I_{11}^{ii(a)}=2\int_{\mathcal{M}\times\mathcal{N}^{\ast}}\beta_{11}^{ii}D_{i}\gamma_{i}A_{i}(|D_{i}z|^{2})-\frac{h^{2}}{2}\int_{\mathcal{M}\times\mathcal{N}^{\ast}}\beta_{11}^{ii}D_{i}\gamma_{i}|D_{i}^{2}z|^{2}.
\end{equation*}
    Then, applying \cite[Lemma 2.2, equation (16)]{LDOP-2021} on the first integral above yields
\begin{equation}
    I_{11}^{ii(a)}=2\int_{\mathcal{M}^{\ast}_{i}\times\mathcal{N}^{\ast}}A_{i}(\beta_{11}^{ii}D_{i}\gamma_{i})|D_{i}z|^{2}+h_{i}\int_{\partial_{i}\mathcal{M}\times\mathcal{N}^{\ast}}\beta_{11}^{ii}D_{i}\gamma_{i}t_{r}^{i}(|D_{i}z|^{2})-\frac{h_{i}^{2}}{2}\int_{\mathcal{M}\times\mathcal{N}^{\ast}}\beta_{11}^{ii}D_{i}\gamma_{i}|D^{2}_{i}z|^{2}.
\end{equation}
 Thus, by using Theorem \ref{theo:weight:estimates}, for $I_{11}^{ii(a)}$ we have 
\begin{equation}\label{eq:I11:ii:a:final}
\begin{aligned}
I_{11}^{ii(a)}=&\;-2\lambda\int_{\mathcal{M}^{\ast}_{i}\times\mathcal{N}^{\ast}}s
\phi\,\partial_{i}\psi\gamma_{i}\,\partial_{i}\gamma_{i}\,|D_{i}z|^{2}
+
\int_{\mathcal{M}^{\ast}_{i}\times\mathcal{N}^{\ast}}s\,\mathcal{O}_{\lambda}((sh)^{2})|D_{i}z|^{2}\\
&+
\int_{\partial_{i}\mathcal{M}\times\mathcal{N}^{\ast}}
\mathcal{O}_{\lambda}(sh)t_{r}^{i}(|D_{i}z|^{2})
+h\,
\int_{\mathcal{M}\times\mathcal{N}^{\ast}}\mathcal{O}_{\lambda}(sh)|D_{i}^{2}z|^{2}.
\end{aligned}
\end{equation}
 Hence, collecting \eqref{eq:J11b} and \eqref{eq:I11:ii:a:final} we get
\begin{equation}
    \begin{aligned}
I_{11}^{ii}=&\int_{\mathcal{M}^{\ast}_{i}\times\mathcal{N}^{\ast}}\left(s\lambda^{2}\gamma_{i}^{2}\phi(\partial_{i}\psi)^{2}-2\lambda s
\phi\,\partial_{i}\psi\gamma_{i}\,\partial_{i}\gamma_{i}+s\lambda\mathcal{O}_{\lambda}(1)+s\mathcal{O}_{\lambda}((sh)^{2})\right)|D_{i}z|^{2}\\
&+\int_{\partial_{i}\mathcal{M}\times\mathcal{N}^{\ast}}(-    s\lambda\gamma_{i}^{2}\phi\partial_{i}\psi+s\mathcal{O}_{\lambda}((sh)^{2}))t_{r}^{i}(|D_{i}z|^{2})n_{i}+h\int_{\mathcal{M}\times\mathcal{N}^{\ast}}\mathcal{O}_{\lambda}(sh)|D_{i}^{2}z|^{2}.
\end{aligned}
\end{equation}
\par Let us focus in the case $i\ne j$. Applying the discrete integration by parts \cite[Lemma 2.2, equation (15)]{LDOP-2021}, we obtain 
\begin{equation*}  I_{11}^{ij}=-2\int_{\mathcal{M}_{i}^{\ast}\times\mathcal{N}^{\ast}}D_{i}( \beta_{11}^{ij}\,D_{j}A_{j}z)\gamma_{i}D_{i}z,
\end{equation*}
since $z=0$ on $\partial\mathcal{M}$ implies $D_{j}A_{j}z=0$ on $\partial_{i}\mathcal{M}$. Now, by equation \eqref{D:product} we write
\begin{equation}
\begin{split}
    I_{11}^{ij}=&-2\int_{\mathcal{M}_{i}^{\ast}\times\mathcal{N}^{\ast}}\gamma_{i}D_{i}( \beta_{11}^{ij})\,D_{j}A_{ij}^{2}z\,D_{i}z-2\int_{\mathcal{M}_{i}^{\ast}\times\mathcal{N}^{\ast}}\gamma_{i}A_{i}(\beta_{11}^{ij})\,D_{ij}^{2}A_{j}z\,D_{i}z=:-I_{11}^{ij(a)}-I_{11}^{ij(b)}.
    \end{split}
\end{equation}
We can integrate by parts with respect to the average operator $A_{i}$ on $I_{11}^{ij(a)}$(\cite[Lemma 2.2, equation (16)]{LDOP-2021}) to obtain
\begin{equation}\label{eq:non:diagonal:I_{11a}}
\begin{split}
    I_{11}^{ij(a)}=&2\int_{\mathcal{M}\times\mathcal{N}^{\ast}}A_{i}(\gamma_{i}D_{i} \beta_{11}^{ij}\,D_{i}z)\,D_{j}A_{j}z\\
    =&2\int_{\mathcal{M}\times\mathcal{N}^{\ast}}A_{i}(\gamma_{i}D_{i} \beta_{11}^{ij})D_{i}A_{i}z\,D_{j}A_{j}z+\frac{h^{2}}{2}\int_{\mathcal{M}\times\mathcal{N}^{\ast}}D_{i}(\gamma_{i}D_{i}\beta_{11}^{ij})D_{i}^{2}zD_{j}A_{j}z.
\end{split}
\end{equation}
where we have used that $D_{j}A_{j}z=0$ on $\partial_{i}\mathcal{M}$ and the identity \eqref{A:product}.
\par Similarly, for $I_{11}^{ij(b)}$, a discrete integration by parts with respect to the average operator $A_{j}$ and product rule given by \eqref{D:product} yield
\begin{equation*}
\begin{split}
    I_{11}^{ij(b)}=&2\int_{\overline{\mathcal{M}}_{ij}\times\mathcal{N}^{\ast}}A_{j}(\gamma_{i}A_{i}\beta_{11}^{ij}\,D_{i}z)\,D_{ij}^{2}z\\
    =&2\int_{\overline{\mathcal{M}}_{ij}\times\mathcal{N}^{\ast}}A_{j}(\gamma_{i}A_{i}\beta_{11}^{ij})\,A_{j}D_{i}z\,D_{ij}^{2}z+\frac{h^{2}}{2}\int_{\overline{\mathcal{M}}_{ij}\times\mathcal{N}^{\ast}}D_{j}(\gamma_{i}A_{i}\beta_{11}^{ij})\,|D_{ij}^{2}z|^{2}.
    \end{split}
\end{equation*}
Now, by virtue of \eqref{D:product} we have $D_{j}\left( |D_{i}z|^{2}\right)=2A_{j}D_{i}z\,D_{ij}^{2}z$, then we can apply a discrete integration by parts with respect to the difference operator $D_{j}$ on the first term from $I_{11}^{ij(b)}$ to get
\begin{equation}\label{eq:non:diagonal:I_{11b}}
\begin{split}
I_{11}^{ij(b)}=&\int_{\overline{\mathcal{M}}_{ij}\times\mathcal{N}^{\ast}}A_{j}(\gamma_{i}A_{i}\beta_{11}^{ij})\,D_{j}\left( |D_{i}z|^{2}\right)+\frac{h^{2}}{2}\int_{\overline{\mathcal{M}}_{ij}\times\mathcal{N}^{\ast}}D_{j}(\gamma_{i}A_{i}\beta_{11}^{ij})\,|D_{ij}^{2}z|^{2}\\
=&-\int_{\mathcal{M}^{\ast}_{i}\times\mathcal{N}^{\ast}}D_{j}A_{j}(\gamma_{i}A_{i}\beta_{11}^{ij})\, |D_{i}z|^{2}+\frac{h^{2}}{2}\int_{\overline{\mathcal{M}}_{ij}\times\mathcal{N}^{\ast}}D_{j}(\gamma_{i}A_{i}\beta_{11}^{ij})\,|D_{ij}^{2}z|^{2},
    \end{split}
\end{equation}
since $D_{i}u=0$ on $\partial_{j}\mathcal{M}^{\ast}_{i}$. Thus, combining \eqref{eq:non:diagonal:I_{11a}} and \eqref{eq:non:diagonal:I_{11b}} we obtain that $I_{11}^{ij}$ can be written as
\begin{equation}
    \begin{split}
       I_{11}^{ij}=&-2\int_{\mathcal{M}\times\mathcal{N}^{\ast}}A_{i}(\gamma_{i}D_{i} \beta_{11}^{ij})D_{i}A_{i}z\,D_{j}A_{j}z-\frac{h^{2}}{2}\int_{\mathcal{M}\times\mathcal{N}^{\ast}}D_{i}(\gamma_{i}D_{i}\beta_{11}^{ij})D_{i}^{2}zD_{j}A_{j}z\\
      &+\int_{\mathcal{M}^{\ast}_{i}\times\mathcal{N}^{\ast}}D_{j}A_{j}(\gamma_{i}A_{i}\beta_{11}^{ij})\, |D_{i}z|^{2}-\frac{h^{2}}{2}\int_{\overline{\mathcal{M}}_{ij}\times\mathcal{N}^{\ast}}D_{j}(\gamma_{i}A_{i}\beta_{11}^{ij})\,|D_{ij}^{2}z|^{2}.
    \end{split}
\end{equation}
Recalling that $\beta_{11}^{ij}:=r^{2}A_{i}^{2}\rho\,D_{j}A_{j}\rho$, Theorem \ref{theo:weight:estimates} and Corollary \ref{Cor:derivative:wrt:x:2r} yield the following estimates
\begin{equation}
    \begin{split}
A_{i}(\gamma_{i}D_{i}\beta_{11}^{ij})&=-s\lambda^{2}\gamma_{i}\gamma_{j}\phi\,\partial_{i}\psi\,\partial_{j}\psi+s\lambda\gamma_{j}\phi\,\mathcal{O}(1)+s\mathcal{O}_{\lambda}((sh)^{2}),\\
        D_{i}(\gamma_{i}D_{i}\beta_{11}^{ij})&=s\mathcal{O}_{\lambda}(1),\\
        D_{j}A_{j}(\gamma_{i}A_{i}\beta_{11}^{ij})&=-s\lambda^{2}\gamma_{i}\gamma_{j}\phi(\partial_{j}\psi)^{2}+s\lambda\gamma_{j}\phi\,\mathcal{O}(1)+s\mathcal{O}_{\lambda}((sh)^{2}),\\
        D_{j}(\gamma_{i}A_{i}\beta_{11}^{ij})&=-s\lambda^{2}\gamma_{i}\gamma_{j}\phi(\partial_{j}\psi)^{2}+s\lambda\gamma_{j}\phi\,\mathcal{O}(1)+s\mathcal{O}_{\lambda}((sh)^{2}).
    \end{split}
\end{equation}
Therefore, by using the estimates above we get
\begin{equation}
    \begin{split}
        I_{11}^{ij}=-\int_{\mathcal{M}_{i}^{\ast}\times\mathcal{N}^{\ast}}s\lambda^{2}\gamma_{i}\gamma_{j}\phi(\partial_{j}\psi)^{2}|D_{i}z|^{2}+V_{11}^{ij}+W_{11}^{ij},
    \end{split}
\end{equation}
where
\begin{equation}
\begin{split}
    W_{11}^{ij}:=&\int_{\mathcal{M}\times\mathcal{N}^{\ast}}\left(s\lambda\varphi\mathcal{O}(1)+s\mathcal{O}_{\lambda}((sh)^{2})\right)\,A_{i}D_{i}z\,D_{j}A_{j}z+\int_{\mathcal{M}\times\mathcal{N}^{\ast}}h\mathcal{O}_{\lambda}(sh)\,D_{i}^{2}z\,D_{j}A_{j}z\\
    &+\int_{\overline{\mathcal{M}}_{ij}\times\mathcal{N}^{\ast}}\left(sh^{2}\lambda\phi\mathcal{O}(1)+sh^{2}\mathcal{O}_{\lambda}((sh)^{2})\right)\,|D_{ij}^{2}z|^{2}\\
    &+\int_{\mathcal{M}_{i}^{\ast}\times\mathcal{N}^{\ast}}(s\lambda\phi\mathcal{O}(1)+s\mathcal{O}_{\lambda}((sh)^{2}))\, |D_{i}z|^{2},
    \end{split}
\end{equation}
\begin{equation}
    V_{11}^{ij}:=\int_{\mathcal{M}\times\mathcal{N}^{\ast}}2s\lambda^{2}\gamma_{i}\gamma_{j}\phi\partial_{j}\psi\partial_{i}\psi\,A_{i}D_{i}z\,D_{j}A_{j}z+\frac{1}{2}\int_{\overline{\mathcal{M}}_{ij}\times\mathcal{N}^{\ast}} sh^{2}\gamma_{i}\gamma_{j}\lambda^{2}\phi(\partial_{j}\psi)^{2}\,|D_{ij}^{2}z|^{2},
\end{equation}

\par On the other hand, we note that 
\begin{equation}
    \begin{split}
        \int_{\mathcal{M}\times\mathcal{N}^{\ast}}h\mathcal{O}_{\lambda}(sh)\,D_{i}^{2}z\,D_{j}A_{j}z\leq &\int_{\mathcal{M}\times\mathcal{N}^{\ast}}h^{2}\mathcal{O}_{\lambda}(sh)\,|D_{i}^{2}z|^{2}+\int_{\mathcal{M}\times\mathcal{N}^{\ast}}\mathcal{O}_{\lambda}(sh)\,|D_{j}A_{j}z|^{2}.
    \end{split}
\end{equation}
Thus, $W_{11}^{ij}$ can be bounded as follows
\begin{equation}
\begin{split}
    |W_{11}^{ij}|:=&\int_{\mathcal{M}\times\mathcal{N}^{\ast}}\left(s\lambda\varphi\mathcal{O}(1)+\mathcal{O}_{\lambda}(sh)+s\mathcal{O}_{\lambda}((sh)^{2})\right)\,|D_{i}A_{i}z|^{2}+\int_{\mathcal{M}\times\mathcal{N}^{\ast}}h^{2}\mathcal{O}_{\lambda}(sh)\,|D_{i}^{2}z|^{2}\\
    &+\int_{\overline{\mathcal{M}}_{ij}\times\mathcal{N}^{\ast}}\left(sh^{2}\lambda\phi\mathcal{O}(1)+sh^{2}\mathcal{O}_{\lambda}((sh)^{2})\right)\,|D_{ij}^{2}z|^{2}\\
    &+\int_{\mathcal{M}^{\ast}_{i}\times\mathcal{N}^{\ast}}(s\lambda\phi\mathcal{O}(1)+s\mathcal{O}_{\lambda}((sh)^{2}))\, |D_{i}z|^{2}\\
    \leq &\int_{\mathcal{M}_{i}^{\ast}\times\mathcal{N}^{\ast}}\left(s\lambda\phi\mathcal{O}(1)+\mathcal{O}_{\lambda}(sh)+s\mathcal{O}_{\lambda}((sh)^{2})\right)\,|D_{i}z|^{2}-\int_{\partial_{i}\mathcal{M}\times\mathcal{N}^{\ast}}\mathcal{O}_{\lambda}(sh)t_{r}^{i}(|D_{i}z|^{2})n_{i}\\
    &+\int_{\mathcal{M}\times\mathcal{N}^{\ast}}h^{2}\mathcal{O}_{\lambda}(sh)\,|D_{i}^{2}z|^{2}+\int_{\overline{\mathcal{M}}_{ij}\times\mathcal{N}^{\ast}}\left(sh^{2}\lambda\varphi\mathcal{O}(1)+sh^{2}\mathcal{O}_{\lambda}((sh)^{2})\right)\,|D_{ij}^{2}z|^{2}\\
    :=&X_{11}^{ij}.
    \end{split}
\end{equation}
From $I_{11}^{ii}$, the leading coefficient of $|D_{i}z|^{2}$ is $s\lambda^{2}\gamma_{i}^{2}\phi(\partial_{i}\psi)^{2}$. From $I_{11}^{ij}$ for each $j\ne i$, the leading coefficient of $|D_{i}z|^{2}$ is $
    -s\lambda^{2}\gamma_{i}\gamma_{j}\phi(\partial_{j}\psi)^{2}$.
Summing over $j\ne i$ we get 
$\displaystyle
    -s\lambda^{2}\gamma_{i}\phi\sum_{j\ne i}\gamma_{j}(\partial_{j}\psi)^{2}$. Then, $\displaystyle s\lambda^{2}\gamma_{i}^{2}\phi(\partial_{i}\psi)^{2}-s\lambda^{2}\gamma_{i}\phi\sum_{j\ne i}\gamma_{j}(\partial_{j}\psi)^{2}=s\lambda^{2}\gamma_{i}\phi\bigg[2\gamma_{i}(\partial_{i}\psi)^{2}-\sum_{j=1}^{d}\gamma_{j}(\partial_{j}\psi)^{2}\bigg].$
Thus, using the notation $|\nabla_{\Gamma}\psi|^{2}:=\sum_{j=1}^{d}\gamma_{j}(\partial_{j}\psi)^{2}$, the final estimate for $I_{11}$ is
\begin{equation}\label{eq:I11:final}
\begin{aligned}
I_{11}=&\sum_{i=1}^{d}\int_{\mathcal{M}_{i}^{\ast}\times\mathcal{N}^{\ast}}s\lambda^{2}\gamma_{i}\phi\Big[2\gamma_{i}(\partial_{i}\psi)^{2}-|\nabla_{\Gamma}\psi|^{2}\Big]|D_{i}z|^{2}\\
&-2s\lambda\sum_{i=1}^{d}\int_{\mathcal{M}_{i}^{\ast}\times\mathcal{N}^{\ast}}\gamma_{i}\phi\,\partial_{i}\psi\,\partial_{i}\gamma_{i}\,|D_{i}z|^{2}+\sum_{i\ne j}V_{11}^{ij}+W_{11}+Y_{11},
\end{aligned}
\end{equation}
where  $V_{11}^{ij}$ is given by
\begin{equation}\label{eq:V11}
V_{11}^{ij}:=2\int_{\mathcal{M}\times\mathcal{N}^{\ast}}s\lambda^{2}\gamma_{i}\gamma_{j}\phi\,\partial_{i}\psi\,\partial_{j}\psi\,D_{i}A_{i}z\,D_{j}A_{j}z+\frac{h^{2}}{2}\int_{\overline{\mathcal{M}}_{ij}\times\mathcal{N}^{\ast}}s\lambda^{2}\gamma_{i}\gamma_{j}\phi(\partial_{j}\psi)^{2}\,|D_{ij}^{2}z|^{2},
\end{equation}
and 
\begin{equation}\label{eq:W11}
\begin{aligned}
W_{11}:=&\sum_{i=1}^{d}\int_{\mathcal{M}_{i}^{\ast}\times\mathcal{N}^{\ast}}\Big(s\lambda\,\mathcal{O}_{\lambda}(1)+s\,\mathcal{O}_{\lambda}((sh)^{2})\Big)|D_{i}z|^{2}+\sum_{i=1}^{d}h\int_{\mathcal{M}\times\mathcal{N}^{\ast}}\mathcal{O}_{\lambda}(sh)\,|D_{i}^{2}z|^{2}\\
&+\sum_{i\ne j}\int_{\mathcal{M}\times\mathcal{N}^{\ast}}\Big(s\lambda\gamma_{j}\phi\,\mathcal{O}(1)+s\,\mathcal{O}_{\lambda}((sh)^{2})\Big)D_{i}A_{i}z\,D_{j}A_{j}z\\
&+\sum_{i\ne j}\int_{\mathcal{M}\times\mathcal{N}^{\ast}}h\,\mathcal{O}_{\lambda}(sh)\,D_{i}^{2}z\,D_{j}A_{j}z\\
&+\sum_{i\ne j}\int_{\overline{\mathcal{M}}_{ij}\times\mathcal{N}^{\ast}}\Big(sh^{2}\lambda\gamma_{j}\phi\,\mathcal{O}(1)+sh^{2}\mathcal{O}_{\lambda}((sh)^{2})\Big)|D_{ij}^{2}z|^{2},
\end{aligned}
\end{equation}
and the boundary terms are
\begin{equation}\label{eq:Y11}
Y_{11}:=\sum_{i=1}^{d}\int_{\partial_{i}\mathcal{M}\times\mathcal{N}^{\ast}}\Big(-s\lambda\gamma_{i}^{2}\phi\,\partial_{i}\psi+\mathcal{O}_{\lambda}(sh)\Big)\,\mathrm{tr}^{i}(|D_{i}z|^{2})\,n_{i}.
\end{equation}

\begin{remark}
\begin{itemize}
    \item
    The coefficient of $|D_{i}z|^{2}$ after summing over $j$ is
    \begin{equation*}
        s\lambda^{2}\gamma_{i}\phi\bigg[2\gamma_{i}(\partial_{i}\psi)^{2}
        -\sum_{j=1}^{d}\gamma_{j}(\partial_{j}\psi)^{2}\bigg].
    \end{equation*}
    When $\gamma_{i}=\gamma$ for all $i$, this equals
    $s\lambda^{2}\gamma^{2}\phi[2(\partial_{i}\psi)^{2}-|\nabla\psi|^{2}]$.
    In dimension $d=1$ this is positive ($+s\lambda^{2}\gamma^{2}\phi|\psi'|^{2}$),
    but for $d\geq 2$ it is not sign-definite. In the full Carleman estimate,
    this term is combined with $I_{21}+I_{22}$ (which provides
    the dominant $s^{3}\lambda^{4}|\nabla\psi|^{4}|z|^{2}$ term)
    and absorbed for $s$ large enough.
    \item The boundary term
    $Y_{11}$ has leading contribution
    $-s\lambda\gamma_{i}^{2}\phi(\partial_{i}\psi\cdot n_{i})\,
    \mathrm{tr}^{i}(|D_{i}z|^{2})$,
    which is positive under the geometric condition
    $\partial_{i}\psi\cdot n_{i}<0$ on $\partial_{i}\mathcal{M}$,
    for $\lambda$ large enough.
\end{itemize}
\end{remark}

%%%%%%%%%%%%%
%%%%%%%%%%%%%%%%%%%%%%%%%%%%%%%%%%%%%%%%
\subsection{Estimate of $I_{12}$ }
%%%%%%%%%%%%%%%%%%%%%%%%%%%%%%%%%%%%%%%%
We set $\beta_{12}^{i}:=-2s\,rA_{i}^{2}\rho\,\Delta_{\Gamma}\phi$. Recall that 
\begin{equation*}
    \begin{split}
        A_{1}z:=&\sum_{i=1}^{d}\trm(rA_{i}^{2}\rho\,D_{i}(\gamma_{i}D_{i}z))\,\text{ and } \,
        B_{2}z:=-2\trm(s\Delta_{\Gamma}\phi z).
    \end{split}
\end{equation*}
Thanks to \cite[Proposition 2, equation (2.4)]{LMPZ:2023}, we can shift the discrete 
time variable to obtain
\begin{equation*}
I_{12}=\sum_{i=1}^{d}\int_{\mathcal{M}\times\mathcal{N}^{\ast}}\beta_{12}^{i}\,z\,D_{i}(\gamma_{i}D_{i}z)=:\sum_{i=1}^{d}I^{i}_{12}.
\end{equation*}
An integration by parts with respect to the difference operator $D_{i}$ yields \break $\displaystyle 
    I_{12}^{i}=-\int_{\mathcal{M}^{\ast}_{i}\times\mathcal{N}^{\ast}}\gamma_{i}D_{i}( \beta_{12}^{i}\,z)\,D_{i}z$,
where we have used that $z=0$ on $\partial_{i}\mathcal{M}$. Now, by using  we write 
    \begin{equation*}
    \begin{split}
    I^{i}_{12}&=-\int_{\mathcal{M}^{\ast}_{i}\times\mathcal{N}^{\ast}}\gamma_{i}D_{i}\beta_{12}^{i}\,A_{i}z\,D_{i}z-\int_{\mathcal{M}^{\ast}_{i}\times\mathcal{N}^{\ast}}\gamma_{i}A_{i}\beta_{12}^{i}\,|D_{i}z|^{2}.
    \end{split}
\end{equation*}
Then, noting that $D_{i}(|z|^{2})=2A_{i}z\,D_{i}z$, $I_{12}^{i}$ can be written as follows
\begin{equation*}
    \begin{split}
    I_{12}^{i}&=-\frac{1}{2}\int_{\mathcal{M}^{\ast}_{i}\times\mathcal{N}^{\ast}}\gamma_{i}D_{i}\beta_{12}^{i}\,D_{i}(|z|^{2})-\int_{\mathcal{M}^{\ast}_{i}\times\mathcal{N}^{\ast}}\gamma_{i}A_{i}(\beta_{12}^{i})|D_{i}z|^{2}.
    \end{split}
\end{equation*}
An integration by parts with respect to the difference operator $D_{i}$(\cite[Lemma 2.2 equation (15)]{LDOP-2021}) on the first integral above gives
\begin{equation}
    \begin{split}
    I_{12}^{i}&=\frac{1}{2}\int_{\mathcal{M}\times\mathcal{N}^{\ast}}D_{i}(\gamma_{i}D_{i}\beta_{12}^{i})\,|z|^{2}-\int_{\mathcal{M}^{\ast}_{i}\times\mathcal{N}^{\ast}}\gamma_{i}A_{i}(\beta_{12}^{i})|D_{i}z|^{2}.
    \end{split}
\end{equation}
By Theorem \ref{theo:weight:estimates} we have $rA_{i}^{2}\rho=1+\mathcal{O}_{\lambda}((sh)^{2})$, then  $\beta_{12}^{ij}=-2s\lambda^{2}\phi|\nabla_{\Gamma}\psi|^{2}-2s\lambda\phi\Delta_{\Gamma} \psi+s\mathcal{O}_{\lambda}((sh)^{2})$ where we have used that $\Delta_{\Gamma}\phi=\lambda\phi(\lambda|\nabla_{\Gamma}\psi|^{2}+\Delta_{\Gamma}\psi)$. 
Moreover,
\begin{equation}\label{eq:beta12:derived}
    \begin{split}
        A_{i}(\beta_{12}^{i})&=-2s\lambda^{2}\phi\,|\nabla_{\Gamma}\psi|^{2}
        -2s\lambda\phi\,\Delta_{\Gamma}\psi+s\,\mathcal{O}_{\lambda}(h^{2}+(sh)^{2}),\\
        D_{i}(\beta_{12}^{i})&=s\,\mathcal{O}_{\lambda}(1),\\
        D_{i}(\gamma_{i}D_{i}\beta_{12}^{i})&=s\,\mathcal{O}_{\lambda}(1).
    \end{split}
\end{equation}
thanks to \cite[Lemma 2.1]{LDOP-2021} and Theorem \ref{theo:weight:estimates}. Thus, $I_{12}$ can be estimate as
\begin{equation}\label{eq:I12:final}
    I_{12}=2\sum_{i=1}^{d}\int_{\mathcal{M}_{i}^{\ast}\times\mathcal{N}^{\ast}}s\lambda^{2}
    \gamma_{i}\phi\,|\nabla_{\Gamma}\psi|^{2}\,|D_{i}z|^{2}+X_{12},
\end{equation}
where the remainder is given by
\begin{equation}\label{eq:X12}
    X_{12}:=\int_{\mathcal{M}\times\mathcal{N}^{\ast}}s\,\mathcal{O}_{\lambda}(1)\,|z|^{2}
    +\sum_{i=1}^{d}\int_{\mathcal{M}_{i}^{\ast}\times\mathcal{N}^{\ast}}
    \Big(2s\lambda\gamma_{i}\phi\,\Delta_{\Gamma}\psi
    +s\,\mathcal{O}_{\lambda}(h^{2}+(sh)^{2})\Big)|D_{i}z|^{2}.
\end{equation}

\begin{remark}
The leading term of $I_{12}$ is
$2s\lambda^{2}\sum_{i}\gamma_{i}\phi\,|\nabla_{\Gamma}\psi|^{2}|D_{i}z|^{2}$,
which is positive and of order $s\lambda^{2}$.
When combined with the leading term of $I_{11}$
(which contains
$s\lambda^{2}\gamma_{i}\phi[2\gamma_{i}(\partial_{i}\psi)^{2}
-|\nabla_{\Gamma}\psi|^{2}]|D_{i}z|^{2}$),
the leading coefficient of $|D_{i}z|^{2}$ becomes
\begin{equation*}
    s\lambda^{2}\gamma_{i}\phi
    \Big[2\gamma_{i}(\partial_{i}\psi)^{2}
    -|\nabla_{\Gamma}\psi|^{2}+2|\nabla_{\Gamma}\psi|^{2}\Big]
    =s\lambda^{2}\gamma_{i}\phi
    \Big[2\gamma_{i}(\partial_{i}\psi)^{2}
    +|\nabla_{\Gamma}\psi|^{2}\Big],
\end{equation*}
which is positive. The sub-leading term $2s\lambda\gamma_{i}\phi\,\Delta_{\Gamma}\psi$
is of order $s\lambda$ and is absorbed for $\lambda$ large.
\end{remark}
%%%%%%%%%%%%
\subsection{Estimate of $I_{13}$}
%%%%%%%%%%%%%%
%%%%%%%%%%%%%%%%%%
For $I_{13}$, let us set $\beta^{i}_{13}:=rA_{i}^{2}\rho$. Then, we have 
\begin{equation*}
I_{13}=\sum_{i=1}^{d}\int_{\mathcal{M}\times\mathcal{N}}\trm(\beta_{13}^{i}\,D_{i}(\gamma_{i}D_{i}z))D_{t}z=:\sum_{i=1}^{d}I^{i}_{13}.
\end{equation*}
Followed by a discrete integration by parts with respect to the difference operator $D_{i}$, $I_{13}$ can be written as
\begin{equation*}
\begin{split}
I^{i}_{13}=&-\int_{\mathcal{M}_{i}^{\ast}\times\mathcal{N}}\trm(\gamma_{i}D_{i}z)\,D_{i}(\trm\beta_{13}^{i}\,D_{t}z),
\end{split}
\end{equation*}
since $D_{t}z=0$ on $\partial\mathcal{M}$. We note that the product rule, \cite[Lemma 2.1]{LDOP-2021}, for the difference operator $D_{i}$ yields
\begin{equation*}
\begin{aligned}
    I_{13}^{i}=&-\int_{\mathcal{M}_{i}^{\ast}\times\mathcal{N}}\trm(\gamma_{i}D_{i}z)D_{i}(\trm\beta_{13}^{i})\,A_{i}D_{t}z\,-\int_{\mathcal{M}_{i}^{\ast}\times\mathcal{N}}\trm(\gamma_{i}D_{i}z)A_{i}\trm(\beta^{i}_{13})\,D_{it}^{2}z\\
    =:&I_{13}^{i(a)}+I_{13}^{i(b)}.
    \end{aligned}
\end{equation*}
Now, by using the identity $\trm zD_{t}z=\frac{1}{2}D_{t}(|z|^{2})-\frac{1}{2}\Delta t|D_{t}z|^{2}$, for $I_{13}^{(b)}$ we get
\begin{equation*}
\begin{split}
    I_{13}^{i(b)}=&-\frac{1}{2}\int_{\mathcal{M}_{i}^{\ast}\times\mathcal{N}}\trm(\gamma_{i}A_{i}\beta_{13}^{i})\,D_{t}(|D_{i}z|^{2})+\frac{\Delta t}{2}\int_{\mathcal{M}_{i}^{\ast}\times\mathcal{N}}\trm(\gamma_{i}A_{i}\beta_{13}^{i})\,|D_{it}^{2}z|^{2}.
    \end{split}
\end{equation*}
Using the integration by parts with respect to the time difference operator $D_{t}$ given by \cite[Proposition 2, equation (2.5)]{LMPZ:2023}, on the first integral of $I_{13}^{i(b)}$, we obtain
\begin{equation*}
\begin{aligned}
    I_{13}^{i(b)}=&\frac{1}{2}\int_{\mathcal{M}_{i}^{\ast}\times\mathcal{N}}\gamma_{i}D_{t}(A_{i}\beta_{13}^{i})\,\trp(|D_{i}z|^{2})-\frac{1}{2}\int_{\mathcal{M}_{i}^{\ast}\times\partial\mathcal{N}}\trp(\gamma_{i}A_{i}\beta_{13}^{i}|D_{i}z|^{2})n\\
    &+\frac{\Delta t}{2}\int_{\mathcal{M}_{i}^{\ast}\times\mathcal{N}}\trm(\gamma_{i}A_{i}\beta_{13}^{i})\,|D_{it}^{2}z|^{2}.
    \end{aligned}
\end{equation*}
By virtue of Theorem \ref{theo:fully:weight:estimates} we note that $A_{i}\taut^{\pm}\alpha_{13}= C_{\lambda}>0$ provided $\tau h(\delta T^{2})^{-1}\leq \varepsilon_{1}(\lambda)$ small enough, we thus have the following lower bound for $I_{13}^{i(b)}$
\begin{equation*}
\begin{split}
    I_{13}^{i(b)}\geq&\frac{1}{2}\int_{\mathcal{M}^{\ast}_{i}\times\mathcal{N}}\gamma_{i}D_{t}(A_{i}\beta_{13}^{i})\,\trp(|D_{i}z|^{2})-C_{\lambda}\int_{\mathcal{M}^{\ast}_{i}\times\partial\mathcal{N}^{+}}\gamma_{i}|D_{i}\trp z|^{2}+C_{\lambda}\Delta t \int_{\mathcal{M}_{i}^{\ast}\times\mathcal{N}}|D_{it}^{2}z|^{2},
    \end{split}
\end{equation*}
where we have dropped the left time boundary term since it is positive 
We note that provided $\frac{\Delta t \tau}{T^{3}\delta^{2}}\leq \frac{1}{2}$, from Theorem \eqref{theo:fully:weight:estimates}, we have 
$\gamma_{i}D_{t}(A_{i}\beta^{i}_{13})=T\trm(\theta(sh)^{2})\mathcal{O}_{\lambda}(1)+\left( \frac{\tau\Delta t}{\delta^{3}T^{4}}\right)\left( \frac{\tau h}{\delta T^{2}}\right)\mathcal{O}_{\lambda}(1)$. Then, we deduce
\begin{equation*}
\begin{split}
    I_{13}^{i(b)}\geq&-\int_{\mathcal{M}^{\ast}_{i}\times\mathcal{N}}\left( T\theta(sh)^{2}\mathcal{O}_{\lambda}(1)+\left( \frac{\tau\Delta t}{\delta^{3}T^{4}}\right)\left( \frac{\tau h}{\delta T^{2}}\right)\mathcal{O}_{\lambda}(1)\right)\,|D_{i}z|^{2}-C_{\lambda}\int_{\mathcal{M}_{i}^{\ast}\times\partial\mathcal{N}^{+}}|D_{i}\trp z|^{2}\\
    &-C_{\lambda}\Delta t\int_{\mathcal{M}_{i}^{\ast}\times\mathcal{N}}\left( \frac{\tau\Delta t}{\delta^{3}T^{4}}\right)\left( \frac{\tau h}{\delta T^{2}}\right)\,|D_{it}^{2}z|^{2}+C_{\lambda}\Delta t \int_{\mathcal{M}_{i}^{\ast}\times\mathcal{N}}|D_{it}^{2}z|^{2},
    \end{split}
\end{equation*}
 where we used $\trp(|D_{i}z|^{2})\leq 2\trm(|D_{i}z|^{2})+2(\Delta t)^{2}|D_{it}z|^{2}$. Hence, taking $\epsilon_{1}(\lambda)$ small enough, for $I_{13}^{i(b)}$ we have
\begin{equation}\label{eq:I13:a2}
\begin{split}
    I_{13}^{i(b)}\geq&-\int_{\mathcal{M}_{i}^{\ast}\times\mathcal{N}}\left( T\theta(sh)^{2}\mathcal{O}_{\lambda}(1)+\left( \frac{\tau\Delta t}{\delta^{3}T^{4}}\right)\left( \frac{\tau h}{\delta T^{2}}\right)\mathcal{O}_{\lambda}(1)\right)\,|D_{i}z|^{2}\\
    &-C_{\lambda}\int_{\mathcal{M}_{i}^{\ast}\times\partial\mathcal{N}^{+}}|D_{i}\trp z|^{2}+C_{\lambda}\Delta t \int_{\mathcal{M}^{\ast}\times\mathcal{N}}|D_{it}^{2}z|^{2}.
    \end{split}
\end{equation}
\par On the other hand, by virtue of Theorem \eqref{theo:weight:estimates}  we note that $D_{i}\beta_{13}^{i}=s\mathcal{O}_{\lambda}((sh)^{2})$. Then, this estimate and the Young's inequality for $I_{13}^{i(a)}$ yield 
\begin{equation*}
    \begin{split}
      |I_{13}^{i(a)}|\leq&C_{\lambda} \int_{\mathcal{M}^{\ast}_{i}\times\mathcal{N}}\trm(s^{-1}(sh)^{2}|D_{t}A_{i}z|^{2}+C_{\lambda}\int_{\mathcal{M}^{\ast}_{i}\times\mathcal{N}^{\ast}}s(sh)^{2}|D_{i}z|^{2},
    \end{split}
\end{equation*}
where in the last integral above we have used \cite[Proposition 2, equation (2.4)]{LMPZ:2023} to shift the discrete time variable. Now, using that $|D_{t}A_{i}z|^{2}\leq A_{i}(|D_{t}z|^{2})$ and an integration by parts with respect to the average operator $A_{i}$ (\cite[Lemma 2.2, equation (16)]{LDOP-2021}) on the first integral of the above expression, we obtain
\begin{equation}\label{eq:I13:a1}
    \begin{split}
      |I_{13}^{i(a)}|\leq&C_{\lambda} \left(\int_{\mathcal{M}\times\mathcal{N}}\trm(s^{-1}(sh)^{2})|D_{t}z|^{2}+\int_{\mathcal{M}_{i}^{\ast}\times\mathcal{N}^{\ast}}s(sh)^{2}|D_{i}z|^{2}\right),
    \end{split}
\end{equation}
since $D_{t}z=0$ on $\partial\mathcal{M}$. Thus, collecting \eqref{eq:I13:a2} and \eqref{eq:I13:a1} we obtain that $I_{13}^{i(a)}$ can be bounded as $\displaystyle I_{13}^{i(a)}\geq -W_{13}$ where
\begin{equation}
    \begin{split}
          W_{13}=&\int_{\mathcal{M}^{\ast}_{i}\times\mathcal{N}^{\ast}}\left( T\theta(sh)^{2}\mathcal{O}_{\lambda}(1)+\left( \frac{\tau\Delta t}{\delta^{3}T^{4}}\right)\left( \frac{\tau h}{\delta T^{2}}\right)\mathcal{O}_{\lambda}(1)\right)\,|D_{i}z|^{2}\\
        &-C_{\lambda}\Delta t \int_{\mathcal{M}_{i}^{\ast}\times\mathcal{N}}|D_{it}^{2}z|^{2}+C_{\lambda}\int_{\mathcal{M}_{i}^{\ast}\times\mathcal{N}^{\ast}}s(sh)^{2}|D_{i}z|^{2}. 
    \end{split}
\end{equation}
Combining the estimates for $I_{13}^{i(a)}$ and $I_{13}^{i(b)}$,
and summing over $i=1,\ldots,d$, we obtain
\begin{equation}\label{eq:I13:final}
\begin{aligned}
I_{13}\geq&-\sum_{i=1}^{d}\int_{\mathcal{M}_{i}^{\ast}\times\mathcal{N}}
\bigg(T\theta(sh)^{2}\,\mathcal{O}_{\lambda}(1)
+\frac{\tau\Delta t}{\delta^{3}T^{4}}\cdot\frac{\tau h}{\delta T^{2}}\,
\mathcal{O}_{\lambda}(1)
+s(sh)^{2}\,\mathcal{O}_{\lambda}(1)\bigg)|D_{i}z|^{2}\\
&-\mathcal{O}_{\lambda}((sh)^{2})
\int_{\mathcal{M}\times\mathcal{N}}|D_{t}z|^{2}
+C_{\lambda}\Delta t\sum_{i=1}^{d}
\int_{\mathcal{M}_{i}^{\ast}\times\mathcal{N}}|D_{it}^{2}z|^{2}\\
&-C_{\lambda}\sum_{i=1}^{d}
\int_{\mathcal{M}_{i}^{\ast}\times\partial\mathcal{N}^{+}}
|D_{i}\tau^{+}z|^{2}.
\end{aligned}
\end{equation}
The positive $\Delta t\int|D_{it}^{2}z|^{2}$ term from $I_{13}^{i(b)}$
is kept on the right-hand side as a good term that can be
moved to the left-hand side of the Carleman estimate.

%%%%%%%%%%%%%%%%%%%%%%%%%%%%%%%%%%%%%
\subsection{Estimate of $I_{21}$}
%%%%%%%%%%%%%%%%%%%%%%%%%%%%%%%%%%%%%
Let us recall that  
\begin{equation*}
\begin{split}
    A_{2}z:=\sum_{i=1}^{d}\gamma_{i}\trm(rD_{i}^{2}\rho A_{i}^{2}z),\, \text{ and }\,
    B_{1}z:=2\sum_{i=1}^{d}\gamma_{i}\trm(rD_{i}A_{i}\rho D_{i}A_{i}z).
    \end{split}
\end{equation*}
Then, setting $\beta_{21}^{ij}:=2\gamma_{i}\gamma_{j}r^{2}D_{i}^{2}\rho D_{j}A_{j}\rho$  we have after shiftting the time varible
\begin{equation*}
    I_{21}=\sum_{i,j=1}^ {d}\int_{\mathcal{M}\times\mathcal{N}^{\ast}}\beta_{21}^{ij}\,A_{i}^{2}z\,D_{j}A_{j}z=:\sum_{i,j=1}^{d}I_{21}^{ij}.
\end{equation*}
\par Let us estimate $I_{21}^{ij}$ for $i\ne j$. Integrating by parts with respect to the average operator $A_{i}$, and using that $D_{j}A_{j}z=0$ on $\partial_{i}\mathcal{M}$, we have 
\begin{equation*}
    \begin{split}
I_{21}^{ij}=\int_{\mathcal{M}_{i}^{\ast}\times\mathcal{N}^{\ast}}A_{i}(\beta_{21}^{ij}\,D_{j}A_{j}z)\,A_{i}z.
    \end{split}
\end{equation*}
Notice that $A_{i}(\beta_{21}^{ij}\,D_{j}A_{j}z)=A_{i}\beta_{21}^{ij}\,A_{ij}^{2}D_{j}z+\frac{h^{2}}{4}D_{i}\beta_{21}^{ij}\,D_{ij}^{2}A_{j}z$, thanks to \cite[Lemma 2.1, equation (14)]{LDOP-2021}. Then, $I_{21}^{ij}$ can be written as
\begin{equation}
    \begin{split}
        I_{21}^{ij}=&\int_{\mathcal{M}_{i}^{\ast}\times\mathcal{N}^{\ast}}A_{i}\beta_{21}^{ij}\,A_{ij}^{2}D_{j}z\,A_{i}z+\frac{h^{2}}{4}\int_{\mathcal{M}_{i}^{\ast}\times\mathcal{N}^{\ast}}D_{i}\beta_{21}^{ij}\,D_{ij}^{2}A_{j}z\,A_{i}z=:I_{21}^{(a)}+I_{21}^{(b)}.
    \end{split}
\end{equation}
\par The next step is to estimate $I_{21}^{(a)}$ and $I_{21}^{(b)}$ as follows. We apply an integration by parts on $I_{21}^{(a)}$ with respect to the average operator $A_{j}$, and use that $A_{i}z=0$ on $\partial_{j}\mathcal{M}$, to obtain $\displaystyle
    I_{21}^{(a)}=\int_{\overline{\mathcal{M}}_{ij}\times\mathcal{N}^{\ast}}A_{j}(A_{i}\beta_{21}^{ij}\,A_{i}z)\,A_{i}D_{j}z$. Then, by using \cite[Lemma 2.1, equation (13)]{LDOP-2021} we have
\begin{equation}
    I_{21}^{(a)}=\frac{1}{2}\int_{\overline{\mathcal{M}}_{ij}\times\mathcal{N}^{\ast}}A^{2}_{ij}\beta_{21}^{ij}\,D_{j}(|A_{i}z|^{2})+\frac{h^{2}}{4}\int_{\overline{\mathcal{M}}_{ij}\times\mathcal{N}^{\ast}}D_{j}A_{i}\beta_{21}^{ij}\,|A_{i}D_{j}z|^{2},
\end{equation}
Thus, applying an integration by parts with respect to the difference operator $D_{j}$ on the first integral in the above expression, we obtain
\begin{equation*}
\begin{split}
    I_{21}^{(a)}=&-\frac{1}{2}\int_{\mathcal{M}^{\ast}_{i}\times\mathcal{N}^{\ast}}D_{j}A^{2}_{ij}\beta_{21}^{ij}\,|A_{i}z|^{2}
    +\frac{h^{2}}{4}\int_{\overline{\mathcal{M}}_{ij}\times\mathcal{N}^{\ast}}D_{j}A_{i}\beta_{21}^{ij}\,|A_{i}D_{j}z|^{2},
    \end{split}
\end{equation*}
where we have used that $A_{i}z=0$ on $\partial_{j}\mathcal{M}$. Moreover, by using $|A_{i}z|^{2}=A_{i}(|z|^{2})-\frac{h^{2}}{4}|D_{i}z|^{2}$ we have
\begin{equation}
\begin{split}
    I_{21}^{(a)}=&-\frac{1}{2}\int_{\mathcal{M}^{\ast}_{i}\times\mathcal{N}^{\ast}}D_{j}A^{2}_{ij}\beta_{21}^{ij}\,A_{i}(|z|^{2})+\frac{h^{2}}{8}\int_{\mathcal{M}^{\ast}_{i}\times\mathcal{N}^{\ast}}D_{j}A^{2}_{ij}\beta_{21}^{ij}\,|D_{i}z|^{2}+\frac{h^{2}}{4}\int_{\overline{\mathcal{M}}_{ij}\times\mathcal{N}^{\ast}}D_{j}A_{i}\beta_{21}^{ij}\,|A_{i}D_{j}z|^{2}.
    \end{split}
\end{equation}
Applying an integration by parts, with respect to the average operator $A_{i}$ on the first term in the above expression, gives
\begin{equation}
\begin{split}
    I_{21}^{(a)}=&-\frac{1}{2}\int_{\mathcal{M}\times\mathcal{N}^{\ast}}A_{i}D_{j}A^{2}_{ij}\beta_{21}^{ij}\,|z|^{2}+\frac{h^{2}}{8}\int_{\mathcal{M}^{\ast}_{i}\times\mathcal{N}^{\ast}}D_{j}A^{2}_{ij}\beta_{21}^{ij}\,|D_{i}z|^{2}+\frac{h^{2}}{4}\int_{\overline{\mathcal{M}}_{ij}\times\mathcal{N}^{\ast}}D_{j}\beta_{21}^{ij}\,|A_{i}D_{j}z|^{2}.
    \end{split}
\end{equation}
Recalling that $\beta_{21}^{ij}:=2\gamma_{i}\gamma_{j}r^{2}D_{i}^{2}\rho D_{j}A_{j}\rho$, from Theorem \ref{theo:weight:estimates} we have
\begin{equation}\label{estimation:beta:21}
    \begin{split}
    \beta_{21}^{ij}=&-2s^{3}\lambda^{3}\phi^{3}\gamma_{i}\gamma_{j}(\partial_{i}\psi)^{2}\partial_{j}\psi+s^{2}\mathcal{O}_{\lambda}(1)+s^{3}\mathcal{O}_{\lambda}((sh)^{2}),\\
        D_{j}\beta_{21}^{ij}=&-6s^{3}\lambda^{4}\phi^{3}\gamma_{i}\gamma_{j}(\partial_{i}\psi)^{2}(\partial_{j}\psi)^{2}+6s^{2}\lambda^{3}\phi^{2}\mathcal{O}_{\lambda}(1)+s^{2}\mathcal{O}_{\lambda}(1)+s^{3}\mathcal{O}_{\lambda}((sh)^{2})=s^{3}\mathcal{O}_{\lambda}(1),\\
        A_{i}D_{j}A_{ij}^{2}\beta_{21}^{ij}=&-6s^{3}\lambda^{4}\phi^{3}(\partial_{i}\psi)^{2}\gamma_{i}\gamma_{j}(\partial_{j}\psi)^{2}+6s^{2}\lambda^{3}\phi^{2}\mathcal{O}_{\lambda}(1)+s^{2}\mathcal{O}_{\lambda}(1)+s^{3}\mathcal{O}_{\lambda}((sh)^{2}).
    \end{split}
\end{equation}
Thus, by using the estimates \eqref{estimation:beta:21}, we get
\begin{equation}\label{eq:J_{21a}}
        I_{21}^{(a)}=\int_{\mathcal{M}\times\mathcal{N}^{\ast}}3s^{3}\lambda^{4}\phi^{3}\gamma_{i}\gamma_{j}(\partial_{i}\psi)^{2}(\partial_{j}\psi)^{2}|z|^{2}+W_{21}^{ij},
\end{equation}
where
\begin{equation}
    \begin{split}
        W_{21}^{ij(a)}:=&\int_{\mathcal{M}\times\mathcal{N}^{\ast}}\left( s^{2}\lambda^{3}\phi^{2}\mathcal{O}_{\lambda}(1)+s^{2}\mathcal{O}_{\lambda}(1)+s^{3}\mathcal{O}_{\lambda}((sh)^{2})\right)|z|^{2}\\
        &+\int_{\mathcal{M}_{i}^{\ast}\times\mathcal{N}^{\ast}}s\mathcal{O}_{\lambda}((sh)^{2})|D_{i}z|^{2}+\int_{\overline{\mathcal{M}}_{ij}\times\mathcal{N}^{\ast}}s\mathcal{O}_{\lambda}((sh)^{2})\,|A_{i}D_{j}z|^{2}.
        \end{split}
\end{equation}
\par Now, we estimate $I_{21}^{(b)}$. By using the estimates \eqref{estimation:beta:21} and applying  Young's inequality we note that
\begin{equation}
\begin{split}
    |I_{21}^{(b)}|\leq&\int_{\mathcal{M}_{i}^{\ast}\times\mathcal{N}^{\ast}}s\mathcal{O}_{\lambda}((sh)^{2})|D_{ij}^{2}A_{j}z|^{2}+\int_{\mathcal{M}_{i}^{\ast}\times\mathcal{N}^{\ast}}s\mathcal{O}_{\lambda}((sh)^{2})|A_{i}z|^{2}.
    \end{split}
\end{equation}
Using inequality $|D_{ij}^{2}A_{j}z|^{2}\leq A_{j}(|D_{ij}^{2}z|^{2})$ it follows that
\begin{equation}
|I_{21}^{(b)}|\leq\int_{\mathcal{M}_{i}^{\ast}\times\mathcal{N}^{\ast}}s\mathcal{O}_{\lambda}((sh)^{2})A_{j}(|D_{ij}^{2}z|^{2})+\int_{\mathcal{M}_{i}^{\ast}\times\mathcal{N}^{\ast}}s\mathcal{O}_{\lambda}((sh)^{2})A_{i}(|z|^{2}).
\end{equation}
Integrating by parts with respect to the average operators $A_{j}$ and $A_{i}$ in the above expression yield
\begin{equation}
|I_{21}^{(b)}|\leq\int_{\overline{\mathcal{M}}_{ij}\times\mathcal{N}^{\ast}}s\mathcal{O}_{\lambda}((sh)^{2})\,|D_{ij}^{2}z|^{2}+\int_{\partial_{j}\mathcal{M}_{i}^{\ast}\times\mathcal{N}^{\ast}}s\mathcal{O}_{\lambda}((sh)^{2})t_{r}^{j}(|D_{ij}^{2}z|^{2})+\int_{\mathcal{M}\times\mathcal{N}^{\ast}}s\mathcal{O}_{\lambda}((sh)^{2})|z|^{2}.
\end{equation}
Hence, when $i\ne j$ we have
\begin{equation}\label{eq:I21:ij:final}
I_{21}^{ij}=3s^{3}\lambda^{4}\int_{\mathcal{M}\times\mathcal{N}^{\ast}}\gamma_{i}\gamma_{j}\phi^{3}(\partial_{i}\psi)^{2}(\partial_{j}\psi)^{2}\,|z|^{2}+W_{21}^{ij},
\end{equation}
where 
\begin{equation}\label{eq:W21:ij:full}
\begin{aligned}
W_{21}^{ij}:=&\int_{\mathcal{M}\times\mathcal{N}^{\ast}}\Big(s^{3}\mathcal{O}_{\lambda}(1)+s^{2}\mathcal{O}_{\lambda}(1)\Big)|z|^{2}+\int_{\mathcal{M}_{i}^{\ast}\times\mathcal{N}^{\ast}}s\,\mathcal{O}_{\lambda}((sh)^{2})|D_{i}z|^{2}\\
&+\int_{\overline{\mathcal{M}}_{ij}\times\mathcal{N}^{\ast}}s\,\mathcal{O}_{\lambda}((sh)^{2})|A_{i}D_{j}z|^{2}+\int_{\overline{\mathcal{M}}_{ij}\times\mathcal{N}^{\ast}}s\,\mathcal{O}_{\lambda}((sh)^{2})|D_{ij}^{2}z|^{2}\\
&+\frac{h}{2}\int_{\partial_{j}\mathcal{M}_{i}^{\ast}\times\mathcal{N}^{\ast}}s\,\mathcal{O}_{\lambda}((sh)^{2})t_{r}^{j}(|D_{ij}^{2}z|^{2})+\frac{h}{2}\int_{\partial_{i}\mathcal{M}\times\mathcal{N}^{\ast}}s\,\mathcal{O}_{\lambda}((sh)^{2})t_{r}^{i}(|z|^{2}).
\end{aligned}
\end{equation}
\par On the other hand, for $j=i$ we have that
\begin{equation*}
I_{21}^{ii}:=\int_{\mathcal{M}\times\mathcal{N}^{\ast}}\beta_{21}^{ii}A_{i}^{2}zD_{i}A_{i}z.
\end{equation*}
Noting that $D_{i}(|A_{i}z|^{2})=2A_{i}^{2}zD_{i}A_{i}z$, a discrete integration by parts gives
\begin{equation*}
I_{21}^{ii}=-\frac{1}{2}\int_{\mathcal{M}^{\ast}_{i}\times\mathcal{N}^{\ast}}D_{i}\beta_{21}^{ii}|A_{i}z|^{2}+\frac{1}{2}\int_{\partial_{i}\mathcal{M}\times\mathcal{N}^{\ast}}\beta_{21}^{ii}t_{r}^{i}(|A_{i}z|^{2})n_{i}.
\end{equation*}
Substituting the identity $|A_{i}z|^{2}=A_{i}(|z|^{2})-\frac{h_{i}^{2}}{4}|D_{i}z|^{2}$ and applying an integration by parts with respect to the average operator $A_{i}$ we obtain
\begin{equation}\label{eq:I21:ii:expanded}
\begin{aligned}
I_{21}^{ii}=&-\frac{1}{2}\int_{\mathcal{M}\times\mathcal{N}^{\ast}}A_{i}(D_{i}\beta_{21}^{ii})\,|z|^{2}+\frac{h^{2}}{8}\int_{\mathcal{M}_{i}^{\ast}\times\mathcal{N}^{\ast}}D_{i}\beta_{21}^{ii}\,|D_{i}z|^{2}\\
&-\frac{h}{4}\int_{\partial_{i}\mathcal{M}\times\mathcal{N}^{\ast}}D_{i}\beta_{21}^{ii}\,t_{r}^{i}(|z|^{2})\,n_{i}+\frac{1}{2}\int_{\partial_{i}\mathcal{M}\times\mathcal{N}^{\ast}}\beta_{21}^{ii}\,t_{r}^{i}(|A_{i}z|^{2})\,n_{i}.
\end{aligned}
\end{equation}
Thanks to Theorem \ref{theo:weight:estimates}
\begin{equation}\label{eq:Di:beta21:ii:estimate}
D_{i}\beta_{21}^{ii}=-6s^{3}\lambda^{4}\gamma_{i}^{2}\phi^{3}(\partial_{i}\psi)^{4}+s^{3}\mathcal{O}_{\lambda}(1)+s^{2}\mathcal{O}_{\lambda}(1)+s^{3}\mathcal{O}_{\lambda}((sh)^{2}).
\end{equation}
and 
\begin{equation}\label{eq:Ai:Di:beta21:ii:estimate}
A_{i}(D_{i}\beta_{21}^{ii})=-6s^{3}\lambda^{4}\gamma_{i}^{2}\phi^{3}(\partial_{i}\psi)^{4}+s^{3}\mathcal{O}_{\lambda}(1)+s^{2}\mathcal{O}_{\lambda}(1)+s^{3}\mathcal{O}_{\lambda}((sh)^{2}).
\end{equation}
Thus, $I_{21}^{ii}$ is estimated as follows
\begin{equation}\label{eq:I21:ii:final}
\begin{aligned}
I_{21}^{ii}=&\;3s^{3}\lambda^{4}\int_{\mathcal{M}\times\mathcal{N}^{\ast}}\gamma_{i}^{2}\phi^{3}(\partial_{i}\psi)^{4}\,|z|^{2}\\
&+\int_{\mathcal{M}\times\mathcal{N}^{\ast}}\Big(s^{3}\mathcal{O}_{\lambda}(1)+s^{2}\mathcal{O}_{\lambda}(1)+s^{3}\mathcal{O}_{\lambda}((sh)^{2})\Big)|z|^{2}\\
&+s\,\mathcal{O}_{\lambda}((sh)^{2})\int_{\mathcal{M}_{i}^{\ast}\times\mathcal{N}^{\ast}}|D_{i}z|^{2}\\
&-\frac{h}{4}\int_{\partial_{i}\mathcal{M}\times\mathcal{N}^{\ast}}\mathcal{O}_{\lambda}(s^{3})\,t_{r}^{i}(|z|^{2})\,n_{i}+\int_{\partial_{i}\mathcal{M}\times\mathcal{N}^{\ast}}s^{3}\mathcal{O}_{\lambda}(1)\,t_{r}^{i}(|A_{i}z|^{2})\,n_{i}.
\end{aligned}
\end{equation}
Therefore, combining \eqref{eq:I21:ij:final} and \eqref{eq:I21:ii:final},
and using the algebraic identity
\begin{equation*}
\sum_{i\ne j}\gamma_{i}\gamma_{j}(\partial_{i}\psi)^{2}(\partial_{j}\psi)^{2}
+\sum_{i=1}^{d}\gamma_{i}^{2}(\partial_{i}\psi)^{4}
=\bigg(\sum_{i=1}^{d}\gamma_{i}(\partial_{i}\psi)^{2}\bigg)^{2},
\end{equation*}
we get
\begin{equation}
I_{21}\geq 3\lambda^{4}\int_{\mathcal{M}\times\mathcal{N}^{\ast}}s^{3}\phi^{3}\left(\sum_{i=1}^{d}\gamma_{i}(\partial_{i}\psi)^{2}\right)^{2}|z|^{2}-X_{21}-Y_{21},
\end{equation}
where
\begin{equation}\label{eq:X21:final}
\begin{aligned}
X_{21}:=&\int_{\mathcal{M}\times\mathcal{N}^{\ast}}
\Big(s^{3}\mathcal{O}_{\lambda}(1)+s^{2}\mathcal{O}_{\lambda}(1)\Big)|z|^{2}
+\sum_{i=1}^{d}
\int_{\mathcal{M}_{i}^{\ast}\times\mathcal{N}^{\ast}}s\,\mathcal{O}_{\lambda}((sh)^{2})|D_{i}z|^{2}\\
&+\sum_{i\ne j}
\int_{\overline{\mathcal{M}}_{ij}\times\mathcal{N}^{\ast}}s\,\mathcal{O}_{\lambda}((sh)^{2})|A_{i}D_{j}z|^{2}
+\sum_{i\ne j})
\int_{\overline{\mathcal{M}}_{ij}\times\mathcal{N}^{\ast}}s\,\mathcal{O}_{\lambda}((sh)^{2}|D_{ij}^{2}z|^{2},
\end{aligned}
\end{equation}
and the boundary remainder is
\begin{equation}\label{eq:Y21:final}
\begin{aligned}
Y_{21}:=&\sum_{i=1}^{d}\bigg[\frac{h}{4}
\int_{\partial_{i}\mathcal{M}\times\mathcal{N}^{\ast}}
s^{3}\mathcal{O}_{\lambda}(1)\,t_{r}^{i}(|z|^{2})\,n_{i}
+\int_{\partial_{i}\mathcal{M}\times\mathcal{N}^{\ast}}
s^{3}\mathcal{O}_{\lambda}(1)\,t_{r}^{i}(|A_{i}z|^{2})\,n_{i}\bigg]\\
&+\sum_{i\ne j}\bigg[\frac{h}{2}
\int_{\partial_{j}\mathcal{M}_{i}^{\ast}\times\mathcal{N}^{\ast}}
s\,\mathcal{O}_{\lambda}((sh)^{2})\,t_{r}^{j}(|D_{ij}^{2}z|^{2})
+\frac{h}{2}\int_{\partial_{i}\mathcal{M}\times\mathcal{N}^{\ast}}
s\,\mathcal{O}_{\lambda}((sh)^{2})\,t_{r}^{i}(|z|^{2})\bigg].
\end{aligned}
\end{equation}
%%%%%%%%%%%%%%%

%%%%%%%%%%%%%%%%%%%%%%%%%%%%%%%%%%%%%
\subsection{Estimate of $I_{22}$ \label{lem:I22}}
%%%%%%%%%%%%%%%%%%%%%%%%%%%%%%%%%%%%%
 We recall that 
\begin{equation}
    A_{2}z:=\sum_{i=1}^{d}\gamma_{i}\trm(rD_{i}^{2}\rho\,A_{i}^{2}z)
\,\text{  and  }\,
    B_{2}z:=-2\trm(s\Delta\phi\,z).
\end{equation}
Let us set $\beta_{22}^{i}:=-2s\gamma_{i}\Delta_{\Gamma} \phi\,rD_{i}^{2}\rho$. Then, shifting the time variable we have
\begin{equation}
I_{22}=\sum_{i=1}^{d}\int_{\mathcal{M}\times\mathcal{N}^{\ast}} \beta_{22}^{i}\,A_{i}^{2}z\,z=:\sum_{i=1}^{d}I_{22}^{i}.
\end{equation}
By using the identity $A_{i}^{2}z=z+\frac{h^{2}}{4}D_{i}^{2}z$ we rewrite $I_{22}^{ij}$ as
\begin{equation}
   I_{22}^{i}=\int_{\mathcal{M}\times\mathcal{N}^{\ast}} \beta_{22}^{i}\,z\,\left(z+\frac{h^{2}}{4}D_{i}^{2}z\right)
    =\int_{\mathcal{M}\times\mathcal{N}^{\ast}} \beta_{22}^{i}|z|^{2}+\frac{h^{2}}{4}\int_{\mathcal{M}\times\mathcal{N}^{\ast}} \beta_{22}^{i}\,z\,D_{i}^{2}z.
\end{equation}
Integration by parts with respect to the difference operator $D_{i}$ yields 
\begin{equation}
    \int_{\mathcal{M}\times\mathcal{N}^{\ast}}\beta_{22}^{i}zD_{i}^{2}z=-\int_{\mathcal{M}_{i}^{\ast}\times\mathcal{N}^{\ast}}D_{i}(\beta_{22}^{i}\,z)\,D_{i}z,
\end{equation}
where we have used that $z=0$ on $\partial_{i}\mathcal{M}$ for $i=1,\ldots,d$. By using product rule in the operator $D_{i}$ \cite[Lemma 2.1, equation (13)]{LDOP-2021} it follows that
\begin{equation}
    I_{22}^{i}=\int_{\mathcal{M}\times\mathcal{N}^{\ast}} \beta_{22}^{i}|z|^{2}-\frac{h^{2}}{4}\int_{\mathcal{M}_{i}^{\ast}\times\mathcal{N}^{\ast}}D_{i}\beta_{22}^{i}\,A_{i}z\,D_{i}z-\frac{h^{2}}{4}\int_{\mathcal{M}_{i}^{\ast}\times\mathcal{N}^{\ast}}A_{i}\beta_{22}^{i}\,|D_{i}z|^{2}.
\end{equation}
Then, from \cite[Lemma 2.1, equation (13)]{LDOP-2021} we rewrite the above expression as
\begin{equation}
I_{22}^{i}=\int_{\mathcal{M}\times\mathcal{N}^{\ast}} \beta_{22}^{i}|z|^{2}-\frac{h^{2}}{8}\int_{\mathcal{M}_{i}^{\ast}\times\mathcal{N}^{\ast}}D_{i}\beta_{22}^{i}\,D_{i}(|z|^{2})-\frac{h^{2}}{4}\int_{\mathcal{M}_{i}^{\ast}\times\mathcal{N}^{\ast}}A_{i}\beta_{22}^{i}\,|D_{i}z|^{2}.
\end{equation}
After integration by parts with respect to the difference operator $D_{i}$, we obtain
\begin{equation*}
    I_{22}^{i}=\int_{\mathcal{M}\times\mathcal{N}^{\ast}} \beta_{22}^{i}|z|^{2}+\frac{h^{2}}{8}\int_{\mathcal{M}\times\mathcal{N}^{\ast}}D_{i}^{2}\beta_{22}^{i}\,|z|^{2}-\frac{h^{2}}{4}\int_{\mathcal{M}_{i}^{\ast}\times\mathcal{N}^{\ast}}A_{i}\beta_{22}^{i}\,|D_{i}z|^{2}.
\end{equation*}
On the other hand, since $\beta_{22}^{i}=-2s\gamma_{i}\Delta_{\Gamma}\phi\cdot rD_{i}^{2}\rho$, with
\begin{equation*}
\Delta_{\Gamma}\phi=\sum_{j=1}^{d}\gamma_{j}\partial_{j}^{2}\phi=\lambda^{2}\phi\sum_{j=1}^{d}\gamma_{j}(\partial_{j}\psi)^{2}+\lambda\phi\sum_{j=1}^{d}\gamma_{j}\partial_{j}^{2}\psi,
\end{equation*}
and, by Theorem~\ref{theo:weight:estimates},
\begin{equation*}
rD_{i}^{2}\rho=s^{2}\lambda^{2}\phi^{2}(\partial_{i}\psi)^{2}+s\,\mathcal{O}_{\lambda}(1)+s^{2}\mathcal{O}_{\lambda}((sh)^{2}),
\end{equation*}
we obtain
\begin{equation}\label{eq:beta22:i:estimate}
\begin{aligned}
\beta_{22}^{i}=&\;-2s^{3}\lambda^{4}\gamma_{i}\phi^{3}(\partial_{i}\psi)^{2}\sum_{j=1}^{d}\gamma_{j}(\partial_{j}\psi)^{2}+s^{3}\lambda^{3}\phi^{3}\,\mathcal{O}(1)+s^{2}\mathcal{O}_{\lambda}(1)+s^{3}\mathcal{O}_{\lambda}((sh)^{2}).
\end{aligned}
\end{equation}
Indeed, the leading term arises from multiplying the dominant parts:
\begin{equation*}
-2s\gamma_{i}\,\lambda^{2}\phi\sum_{j=1}^{d}\gamma_{j}(\partial_{j}\psi)^{2}\, s^{2}\lambda^{2}\phi^{2}(\partial_{i}\psi)^{2}=-2s^{3}\lambda^{4}\gamma_{i}\phi^{3}(\partial_{i}\psi)^{2}\sum_{j=1}^{d}\gamma_{j}(\partial_{j}\psi)^{2},
\end{equation*}
while term $s^{3}\lambda^{3}\phi^{3}\,\mathcal{O}(1)$ comes from the cross products
\begin{equation*}
-2s\gamma_{i}\,\lambda\phi\sum_{j=1}^{d}\gamma_{j}\partial_{j}^{2}\psi\, s^{2}\lambda^{2}\phi^{2}(\partial_{i}\psi)^{2}=-2s^{3}\lambda^{3}\gamma_{i}\phi^{3}(\partial_{i}\psi)^{2}\sum_{j=1}^{d}\gamma_{j}\partial_{j}^{2}\psi=s^{3}\lambda^{3}\phi^{3}\,\mathcal{O}(1).
\end{equation*}
The remaining terms $s^{2}\mathcal{O}_{\lambda}(1)$ and $s^{3}\mathcal{O}_{\lambda}((sh)^{2})$ arise from the lower-order contributions in $rD_{i}^{2}\rho$ and the discrete error in the theorem. Consequently, 
\begin{equation*}
\sum_{i=1}^{d}\beta_{22}^{i}=-2s^{3}\lambda^{4}\phi^{3}\bigg(\sum_{i=1}^{d}\gamma_{i}(\partial_{i}\psi)^{2}\bigg)^{2}+s^{3}\lambda^{3}\phi^{3}\,\mathcal{O}(1)+s^{2}\mathcal{O}_{\lambda}(1)+s^{3}\mathcal{O}_{\lambda}((sh)^{2}).
\end{equation*}
Moreover, we also have
\begin{equation*}
D_{i}^{2}\beta_{22}^{i}=s^{3}\mathcal{O}_{\lambda}(1),\qquad A_{i}\beta_{22}^{i}=s^{3}\mathcal{O}_{\lambda}(1)+s\,\mathcal{O}_{\lambda}((sh)^{2}),
\end{equation*}
Thus, by using the above bounds and \eqref{eq:beta22:i:estimate}, $I_{22}$ can be estimated as 
\begin{equation}\label{eq:I22:final}
I_{22}=-2\lambda^{4}\int_{\mathcal{M}\times\mathcal{N}^{\ast}}s^{3}\phi^{3}\bigg(\sum_{i=1}^{d}\gamma_{i}(\partial_{i}\psi)^{2}\bigg)^{2}|z|^{2}+X_{22},
\end{equation}
where 
\begin{equation}\label{eq:X22}
\begin{aligned}
X_{22}:=&\int_{\mathcal{M}\times\mathcal{N}^{\ast}}\Big(s^{3}\lambda^{3}\phi^{3}\,\mathcal{O}(1)+s^{2}\mathcal{O}_{\lambda}(1)+s^{3}\mathcal{O}_{\lambda}((sh)^{2})\Big)|z|^{2}\\
&+\int_{\mathcal{M}\times\mathcal{N}^{\ast}}s\,\mathcal{O}_{\lambda}((sh)^{2})|z|^{2}+\sum_{i=1}^{d}\int_{\mathcal{M}_{i}^{\ast}\times\mathcal{N}^{\ast}}s\,\mathcal{O}_{\lambda}((sh)^{2})|D_{i}z|^{2}.
\end{aligned}
\end{equation}
\begin{remark}
The leading term of $I_{22}$ is $\displaystyle -2\lambda^{4}\int s^{3}\phi^{3}\bigg(\sum_{i}\gamma_{i}(\partial_{i}\psi)^{2}\bigg)^{2}|z|^{2}$, which is negative. When combined with the positive leading term of $I_{21}$, namely 
$$3s^{3}\lambda^{4}\int\phi^{3}\bigg(\sum_{i}\gamma_{i}(\partial_{i}\psi)^{2}\bigg)^{2}|z|^{2},$$ the contribution is
\begin{equation*}
I_{21}+I_{22}=s^{3}\lambda^{4}\int_{\mathcal{M}\times\mathcal{N}^{\ast}}\phi^{3}\bigg(\sum_{i=1}^{d}\gamma_{i}(\partial_{i}\psi)^{2}\bigg)^{2}|z|^{2}+\widetilde{W}_{21}+X_{22},
\end{equation*}
which provides the main positive $s^{3}\lambda^{4}|\nabla\psi|^{4}|z|^{2}$ term in the Carleman estimate (when $\gamma_{i}=\gamma$ for all $i$, the squared sum becomes $\gamma^{2}|\nabla\psi|^{4}$). The sub-leading terms in $X_{22}$ are absorbed for $\lambda$ large (the $s^{3}\lambda^{3}$ term) and for $sh$ small (the $(sh)^{2}$ terms).
\end{remark}
%%%%%%%%%%%%%%%%%
\subsection{Estimate of $I_{23}$}
%%%%%%%%%%%%
%%%%%%%%%%
Let us recall that
\begin{equation*}
    A_{2}z:=\sum_{i=1}^{d}\gamma_{i}\trm(rD_{i}^{2}\rho A_{i}^{2}z)\, \text{ and }B_{3}z:=D_{t}z.
\end{equation*}
Then, setting $\beta_{23}:=rD_{i}^{2}\rho$, we have to estimate
\begin{equation}
I_{23}=\sum_{i=1}^{d}\int_{\mathcal{M}\times\mathcal{N}}\gamma_{i}\trm(\beta_{23}A_{i}^{2}z)\,D_{t}z=:\sum_{i=1}^{d}I^{i}_{23}.
\end{equation}
Using the identity $A_{i}^{2}z=z+\frac{h^{2}}{4}D_{i}^{2}z$, $I^{i}_{23}$ can be rewritten as
\begin{equation}
\begin{split}
    I^{i}_{23}=&\int_{\mathcal{M}\times\mathcal{N}}\gamma_{i}\trm(\beta_{23})\trm(z)\,D_{t}z+\frac{h^{2}}{4}\int_{\mathcal{M}\times\mathcal{N}}\gamma_{i}\trm(\beta_{23}\,D_{i}^{2}z)\,D_{t}z\\
    =:&I_{23}^{i(a)}+I_{23}^{i(b)}.
    \end{split}
    \end{equation}
\par We note that the identity $\trm z\,D_{t}z=\frac{1}{2}D_{t}(|z|^{2})-\frac{\Delta t}{2}|D_{t}z|^{2}$, for $I_{23}^{i(a)}$, yields 
\begin{equation*}
\begin{split}
    I_{23}^{i(a)}=&\frac{1}{2}\int_{\mathcal{M}\times\mathcal{N}}\gamma_{i}\trm(\beta_{23})\,D_{t}(|z|^{2})-\frac{\Delta t}{2}\int_{\mathcal{M}\times\mathcal{N}}\gamma_{i}\trm(\beta_{23})\,|D_{t}z|^{2}.
    \end{split}
    \end{equation*}
A discrete integration by parts with respect to the time difference operator $D_{t}$, \cite[Equation (2.9) ]{LMPZ:2023} enables us to write
\begin{equation*}
\begin{split}
    I_{23}^{i(a)}=&-\frac{1}{2}\int_{\mathcal{M}\times\mathcal{N}}\gamma_{i}D_{t}(\beta_{23})\trp(|z|^{2})+\frac{1}{2}\int_{\mathcal{M}\times\partial\mathcal{N}}\gamma_{i}\trp(\beta_{23}|z|^{2})\,n\\
    &-\frac{\Delta t}{2}\int_{\mathcal{M}\times\mathcal{N}}\gamma_{i}\trm(\beta_{23})\,|D_{t}z|^{2}.
\end{split}
\end{equation*}
Thus, by using Theorem \ref{theo:fully:weight:estimates} and using the inequality $|\trp(|z|^{2})|\leq 2\trm(|z|^{2})+2(\Delta t)^{2}|D_{t}z|^{2}$, for $I_{23}^{i(a)}$ we obtain
\begin{equation}\label{ine:I:23:a}
    \begin{aligned}
    I_{23}^{i(a)}=&-\int_{\mathcal{M}\times\mathcal{N}}\trm\mu\,\trm(|z|^{2})-(\Delta t)^{2}\int_{\mathcal{M}\times\mathcal{N}}\trm(Ts^{2}\theta)|D_{t}z|^{2}\\
    &+\int_{\mathcal{M}\times\partial\mathcal{N}}\mathcal{O}_{\lambda}(1)\trp(s^{2}|z|^{2})\,n-\Delta t\int_{\mathcal{M}\times\mathcal{N}}\trm(s^{2})\mathcal{O}_{\lambda}(1)\,|D_{t}z|^{2},
    \end{aligned}
\end{equation}
where $\mu:=Ts^{2}\theta\mathcal{O}_{\lambda}(1)+\left(\frac{\tau^{2}\Delta t}{\delta^{4}T^{6}}\right)\mathcal{O}_{\lambda}(1)+\left(\frac{\tau \Delta t}{\delta^{3}T^{4}}\right)\left(\frac{\tau h}{\delta T^{2}}\right)^{2}\mathcal{O}_{\lambda}(1)$.
\par Let us now estimate $I_{23}^{i(b)}$. Using a discrete integral by parts for the operator $D_{i}$, and $D_{t}z=0$ on $\partial\mathcal{M}$, we obtain
\begin{equation*}
    I_{23}^{(b)}=-\frac{h^{2}}{4}\int_{\mathcal{M}_{i}^{\ast}\times\mathcal{N}}D_{i}(\gamma_{i}\trm\beta_{23}D_{t}z)\trm(D_{i}z).
\end{equation*}
By virtue of the product rule for the difference operator $D_{i}$ we write
\begin{equation*}
\begin{split}
    I_{23}^{(b)}=&-\frac{h^{2}}{4}\int_{\mathcal{M}_{i}^{\ast}\times\mathcal{N}}\trm(D_{i}(\gamma_{i}\beta_{23}))A_{i}D_{t}z\,\trm(D_{i}z)-\frac{h^{2}}{4}\int_{\mathcal{M}_{i}^{\ast}\times\mathcal{N}}\trm(A_{i}(\gamma_{i}\beta_{23}))D^{2}_{it}z\,\trm(D_{i}z)\\
    :=&I_{23}^{(b_{1})}+I_{23}^{(b_{2})}.
    \end{split}
\end{equation*}
We note that $D_{it}^{2}z\,\trm(D_{i}z)=\frac{1}{2}D_{t}(|D_{i}z|^{2})-\frac{\Delta t}{2}|D^{2}_{ti}z|^{2}$, thanks to the identity $\trm z\,D_{t}z=\frac{1}{2}D_{t}(|z|^{2})-\frac{\Delta t}{2}|D_{t}z|^{2}$. Then, $I_{23}^{(b_{2})}$ can be rewritten as
\begin{equation}
\begin{aligned}
    I_{23}^{(b_{2})}=&-\frac{h^{2}}{8}\int_{\mathcal{M}_{i}^{\ast}\times\mathcal{N}}\trm(A_{i}(\gamma_{i}\beta_{23}))D_{t}(|D_{i}z|^{2})+\frac{\Delta t\,h^{2}}{8}\int_{\mathcal{M}_{i}^{\ast}\times\mathcal{N}}\trm(A_{i}(\gamma_{i}\beta_{23})) 
    |D^{2}_{ti}z|^{2}.
    \end{aligned}
\end{equation}
Applying a discrete integration by parts for the discrete difference time operator $D_{t}$ on the first integral above leads to
\begin{equation}
\begin{split}
    I_{23}^{(b_{2})}=&\frac{h^{2}}{8}\int_{\mathcal{M}^{\ast}\times\mathcal{N}}A_{i}D_{t}(\gamma_{i}\beta_{23})\trp(|D_{i}z|^{2})-\frac{h^{2}}{8}\int_{\mathcal{M}^{\ast}\times\partial\mathcal{N}}\trp(A_{i}(\gamma_{i}\beta_{23})|D_{i}z|^{2})\,n\\
    &+\frac{\Delta th^{2}}{8}\int_{\mathcal{M}_{i}^{\ast}\times\mathcal{N}}\trm(A_{i}(\gamma_{i}\beta_{23}))\,|D^{2}_{ti}z|^{2}.
    \end{split}
\end{equation}
We note that using the estimates Theorem \ref{theo:fully:weight:estimates} and the estimates for the spatial variable it follows that
\begin{equation}\label{eq:I_{23}:}
    \begin{split}
        A_{i}D_{t}(\gamma_{i}\beta_{23})&=T\trm(s^{2}\theta)\mathcal{O}_{\lambda}(1)+\frac{\tau^{2}\Delta t}{\delta^{4} T^{6}}\mathcal{O}_{\lambda}(1)+\left( \frac{\tau \Delta t}{\delta^{3}T^{4}}\right)\left( \frac{\tau h}{\delta T^{2}}\right)^{2}\mathcal{O}_{\lambda}(1),\\
        A_{i}(\gamma_{i}\beta_{23})&=s^{2}\mathcal{O}_{\lambda}(1).
    \end{split}
\end{equation}
Combining the previous estimates and \eqref{eq:I_{23}:}, and using $\trp(|D_{i}z|^{2})\leq 2\trm(|D_{i}z|^{2})+2(\Delta t)^{2}|D_{ti}^{2}z|^{2}$, we have $I_{23}^{(b_{2})}$ can be estimated as 
\begin{equation}\label{ine:I:23:b2}
\begin{aligned}
|I_{23}^{(b_{2})}|\leq&\int_{\mathcal{M}^{\ast}\times\mathcal{N}}T\trm(\theta (sh)^{2})\trm(|D_{i}z|^{2})+(\Delta t)^{2}\int_{\mathcal{M}^{\ast}\times\mathcal{N}}T\trm(\theta (sh)^{2})|D^{2}_{it}z|^{2})\\
&+h^{2}\int_{\mathcal{M}^{\ast}\times\partial\mathcal{N}}\frac{\tau^{2}\Delta t}{\delta^{4} T^{6}}\mathcal{O}_{\lambda}(1)\trp(|D_{i}z|^{2})+\Delta th^{2}\int_{\mathcal{M}\times\mathcal{N}}\mathcal{O}_{\lambda}(1)\trm(s^{2})\,|D^{2}_{ti}z|^{2}\\
&+h^{2}\int_{\mathcal{M}^{\ast}\times \partial \mathcal{N}}\left( \frac{\tau \Delta t}{\delta^{3}T^{4}}\right)\left( \frac{\tau h}{\delta T^{2}}\right)^{2}\mathcal{O}_{\lambda}(1)\trp(|D_{i}z|^{2}).
    \end{aligned}
\end{equation}
\par On the other hand, we focus now on $I_{23}^{(b_{1})}$. We note that Young's inequality yields
\begin{equation}
    |I_{23}^{(b_{1})}|\leq C \int_{\mathcal{M}^{\ast}\times\mathcal{N}}h^{2}|\trm(s^{-1}D_{i}(\gamma_{i}\beta_{23}))|A_{i}D_{t}z|^{2}+\int_{\mathcal{M}^{\ast}\times\mathcal{N}}h^{2}|\trm(D_{i}(\gamma_{i}\beta_{23}))|\trm(s|D_{i}z|^{2}).
\end{equation}
Using that $|A_{i}D_{t}z|^{2}\leq A_{i}(|D_{t}z|^{2})$ on the first integral from the above expression and then a discrete integration by part for the average operator $A_{i}$ we have
\begin{equation}
\begin{split}
    |I_{23}^{(b_{1})}|\leq &C \int_{\mathcal{M}\times\mathcal{N}}h^{2}\trm(
    s^{-1}|A_{i}D_{i}(\gamma_{i}\beta_{23})|)\,|D_{t}z|^{2}+\int_{\mathcal{M}^{\ast}\times\mathcal{N}}h^{2}|\trm(D_{i}\alpha_{23})|\trm(s|D_{i}z|^{2}),
    \end{split}
\end{equation}
where we have used $D_{t}z=0$ on $\partial_{i}\mathcal{M}$. Finally, using Theorem \ref{theo:fully:weight:estimates}, $I_{23}^{(b_{1})}$ can be estimated as follows
\begin{equation}\label{ine:I:23:b1}
    |I_{23}^{(b_{1})}|\leq C \int_{\mathcal{M}\times\mathcal{N}}s^{-1}\mathcal{O}_{\lambda}((sh)^{2})\,|D_{t}z|^{2}+\int_{\mathcal{M}^{\ast}\times\mathcal{N}}\mathcal{O}_{\lambda}((sh)^{2})\trm(s|D_{i}z|^{2}).
\end{equation}
Thus, combining estimates \eqref{ine:I:23:a}, \eqref{ine:I:23:b1} and \eqref{ine:I:23:b2} we get
\begin{equation}\label{eq:I23:final}
\begin{aligned}
I_{23}^{i}=&-\int_{\mathcal{M}\times\mathcal{N}}\gamma_{i}\mu\,\trm(|z|^{2})+X_{23}^{i}+Y_{23}^{i},
\end{aligned}
\end{equation}
where $\mu:=Ts^{2}\theta\,\mathcal{O}_{\lambda}(1)+\frac{\tau^{2}\Delta t}{\delta^{4}T^{6}}\,\mathcal{O}_{\lambda}(1)+\frac{\tau^{4}h^{2}\Delta t}{\delta^{6}T^{10}}\,\mathcal{O}_{\lambda}(1)$, the volume remainder $\mathcal{R}_{23}^{i}$ satisfies
\begin{equation}\label{eq:R23}
\begin{aligned}
X_{23}^{i}:=&-\bigg[(\Delta t)^{2}Ts^{2}\theta\,\mathcal{O}_{\lambda}(1)+\Delta t\,s^{2}\,\mathcal{O}_{\lambda}(1)+s^{-1}\mathcal{O}_{\lambda}((sh)^{2})\bigg]\int_{\mathcal{M}\times\mathcal{N}}|D_{t}z|^{2}\\
&+\bigg[T\theta(sh)^{2}\,\mathcal{O}_{\lambda}(1)+\mathcal{O}_{\lambda}((sh)^{2})\,s\bigg]\int_{\mathcal{M}_{i}^{\ast}\times\mathcal{N}}\trm(|D_{i}z|^{2})\\
&+(\Delta t)^{2}T\theta(sh)^{2}\,\mathcal{O}_{\lambda}(1)\int_{\mathcal{M}_{i}^{\ast}\times\mathcal{N}}|D_{ti}^{2}z|^{2}\\
&+\Delta t\,h^{2}\,s^{2}\,\mathcal{O}_{\lambda}(1)\int_{\mathcal{M}_{i}^{\ast}\times\mathcal{N}}\trm(|D_{ti}^{2}z|^{2}),
\end{aligned}
\end{equation}
and the boundary remainder $Y_{23}^{i}$ collects the contributions on $\partial\mathcal{N}$:
\begin{equation}\label{eq:B23}
\begin{aligned}
Y_{23}^{i}:=&\;\frac{1}{2}\int_{\mathcal{M}\times\partial\mathcal{N}}\gamma_{i}\,\mathcal{O}_{\lambda}(1)\,\trp(s^{2}|z|^{2})\,n-\frac{h^{2}}{8}\int_{\mathcal{M}_{i}^{\ast}\times\partial\mathcal{N}}\trp(A_{i}(\gamma_{i}\beta_{23})\,|D_{i}z|^{2})\,n.
\end{aligned}
\end{equation}
\par Summing over $i=1,\ldots,d$, the full estimate for $I_{23}$ reads

\begin{equation}\label{eq:I23:total}
\begin{aligned}
I_{23}=&-\sum_{i=1}^{d}\int_{\mathcal{M}\times\mathcal{N}}\gamma_{i}\mu\,\trm(|z|^{2})+\sum_{i=1}^{d}X_{23}^{i}+\sum_{i=1}^{d}Y_{23}^{i}.
\end{aligned}
\end{equation}
\begin{remark}
The terms in $X_{23}^{i}$ have the following roles in the Carleman estimate:
\begin{itemize}
    \item The $|D_{t}z|^{2}$ terms: the leading contribution $\Delta t\,s^{2}\,\mathcal{O}_{\lambda}(1)$ must be absorbed by the Carleman term $\frac{1}{s}\int|D_{t}z|^{2}$. The $(\Delta t)^{2}Ts^{2}\theta$ and $s^{-1}(sh)^{2}$ terms are lower order.
    \item The $\trm(|D_{i}z|^{2})$ terms: the leading contribution $T\theta(sh)^{2}$ is absorbed for $sh$ small, and $s(sh)^{2}$ is lower order.
    \item The $|D_{ti}^{2}z|^{2}$ terms: both carry factors of $(\Delta t)^{2}$ or $\Delta t\,h^{2}$, making them lower order.
\end{itemize}
\end{remark}

%%%%%%%%%%%%%%%
%%%%%%%%%
\subsection{Estimate of $I_{31}$}
%%%%%%%%%
\begin{lemma}\label{I31}
Provided $\Delta t \tau (T^3 \delta^2)^{-1}\leq 1$ and $\tau h(\delta T^{2})^{-1}\leq 1$, we have
\begin{equation*}
\begin{split}
    I_{31}\geq &X_{31}+Y_{31}
    ,
    \end{split}
\end{equation*}
where $\displaystyle X_{31}:=\int_{\mathcal{M}\times\mathcal{N}^{\ast}}\lambda\varphi\partial_{t}\theta s\mathcal{O}_{\lambda}(1)\,|z|^{2}+\int_{\mathcal{M}^{\ast}\times\mathcal{N}^{\ast}}s\mathcal{O}_{\lambda}(1)\,|D_{h}z|^{2}$ and
\begin{equation*}
    Y_{31}:=-h\int_{\partial \mathcal{M}\times\mathcal{N}^{\ast}}\lambda\varphi\partial_{t}\theta s\mathcal{O}_{\lambda}(1)t_{r}(|D_{h}z|^{2}).
\end{equation*}
\end{lemma}
%%%%%%%%%
Recall that 
$A_{3}z=-\tau\varphi\trm(\theta'z)$ and $B_{1}z=2\sum_{i=1}^{d}\gamma_{i}\trm(rD_{i}A_{i}\rho D_{i}A_{i}z)$. Let us set $\alpha_{31}:=-\tau\varphi\theta'\gamma_{i}rD_{i}A_{i}\rho$. Then, shifting the time variable for $I_{31}$ yields $$\displaystyle I_{31}=\sum_{i=1}^{d}\int_{\mathcal{M}\times\mathcal{N}^{\ast}}2\alpha_{31}z\,D_{i}A_{i}z.$$ A discrete integration by parts with respect to the average operator $A_{i}$ leads to
\begin{equation*}
I^{i}_{31}=\int_{\mathcal{M}_{i}^{\ast}\times\mathcal{N}^{\ast}}2A_{i}(\alpha_{31}z)D_{i}z.
\end{equation*}
From \cite[Lemma 2.1, equation (13)]{LDOP-2021} and the identity $2D_{i}zA_{i}z=D_{i}(|z|^{2})$, for $I_{31}^{(a)}$, we have
\begin{equation*}
\begin{split}
    I_{31}^{i}=&\int_{\mathcal{M}_{i}^{\ast}\times\mathcal{N}^{\ast}}A_{i}(\alpha_{31})\,D_{i}(|z|^{2})+\frac{h^{2}}{2}\int_{\mathcal{M}_{i}^{\ast}\times\mathcal{N}^{\ast}}D_{i}(\alpha_{31})\,|D_{i}z|^{2}.
    \end{split}
\end{equation*}
A discrete integration by parts with respect to the difference operator $D_{i}$ yields
\begin{equation*}
\begin{split}
    I_{31}^{(i)}=&-\int_{\mathcal{M}\times\mathcal{N}^{\ast}}D_{i}A_{i}(\alpha_{31})\,|z|^{2}+\frac{h^{2}}{2}\int_{\mathcal{M}_{i}^{\ast}\times\mathcal{N}^{\ast}}D_{i}(\alpha_{31})\,|D_{i}z|^{2},
    \end{split}
\end{equation*}
since $z=0$ on $\partial_{i}\mathcal{M}$ for $i=1,\ldots, d$. We note that $e^{s\varphi}\partial_{t}(e^{-s\varphi})=-\tau\varphi\partial_{t}\theta$ and by Theorem \ref{theo:weight:estimates} we have $rA_{i}D_{i}\rho=s\mathcal{O}_{\lambda}(1)$. Thus, for $I_{31}^{i}$ we have the following estimate
\begin{equation}
\begin{split}
I_{31}^{(i)}=&\int_{\mathcal{M}\times\mathcal{N}^{\ast}}Ts^{2}\theta\mathcal{O}_{\lambda}(1)\,|z|^{2}+\int_{\mathcal{M}^{\ast}\times\mathcal{N}^{\ast}}T\theta\mathcal{O}_{\lambda}((sh)^{2})\,|D_{i}z|^{2}.
    \end{split}
\end{equation}
Therefore, $I_{31}$ can be estimated as 
\begin{equation*}
I_{31}=\int_{\mathcal{M}\times\mathcal{N}^{\ast}}dTs^{2}\theta\mathcal{O}_{\lambda}(1)\,|z|^{2}+\sum_{i=1}^{d}\int_{\mathcal{M}_{i}^{\ast}\times\mathcal{N}^{\ast}}T\theta\mathcal{O}_{\lambda}((sh)^{2})\,|D_{i}z|^{2}.
\end{equation*}
%%%%%%%%%%
%%%%%%%
\subsection{Estimate of $I_{32}$}
%%%%%%
Let us set $\alpha_{32}:=2s\Delta_{\Gamma}\phi\tau\varphi\theta'$. We have shifting the discrete time variable $I_{32}$ is given by $\displaystyle I_{32}=\int_{\mathcal{M}\times\mathcal{N}^{\ast}}\alpha_{32}|z|^{2}$. Then, it follows that 
$$I_{32}\geq -\int_{\mathcal{M}\times\mathcal{N}^{\ast}}Ts^{2}\theta\mathcal{O}_{\lambda}(1)|z|^{2},$$
since $|\theta'|\leq CT\theta^{2}$, $\left\| \Delta_{\Gamma}\phi\right\|_{L^{\infty}_{h}(\mathcal{M}\times\mathcal{N}^{\ast})}\leq \mathcal{O}_{\lambda}(1)$ and $\left\| \varphi\right\|_{C(\Omega)}=\mathcal{O}_{\lambda}(1)$.
%%%%%%%%%%%%
%%%%%%%%
\subsection{Estimate of $I_{33}$}
%%%%%%%%%%%
We have to estimate
$\displaystyle I_{33}:=-\int_{\mathcal{M}\times\mathcal{N}}\tau\varphi\trm(\theta'\,z)D_{t}z$. Note that by using the identity $\trm(z)D_{t}z=\frac{1}{2}D_{t}(|z|^{2})-\frac{\Delta t}{2}|D_{t}z|^{2}$, $I_{33}$ can be rewritten as 
\begin{equation*}
I_{33}=-\frac{1}{2}\int_{\mathcal{M}\times\mathcal{N}}\tau\varphi\trm(\theta')\,D_{t}(|z|^{2})+\frac{\Delta t}{2}\int_{\mathcal{M}\times\mathcal{N}}\tau\varphi\trm(\theta')|D_{t}z|^{2}.
\end{equation*}
Now, using \cite[Equation (2.9)]{LMPZ:2023} on the first integral above we have 
\begin{equation*}
\begin{split}
    I_{33}=&\frac{1}{2}\int_{\mathcal{M}\times\mathcal{N}}\tau\varphi\,D_{t}(\theta')\,\trp(|z|^{2})-\frac{1}{2}\int_{\mathcal{M}\times\partial \mathcal{N}}\tau\varphi\trp(|z|^{2})\trp(\theta')n+\frac{\Delta t}{2}\int_{\mathcal{M}\times\mathcal{N}}\tau\varphi\trm(\theta')|D_{t}z|^{2}.
    \end{split}
\end{equation*}
We note that $\trp(\theta')\,n>0$ on $\mathcal{M}\times\partial\mathcal{N}$ since $\theta'=(2t-T)\theta^{2}$. Moreover, thanks to $\varphi<0$ in $\mathcal{M}$, it follows that
\begin{equation}
-\frac{1}{2}\int_{\mathcal{M}\times\partial\mathcal{N}}\tau\varphi\trp(|z|^{2})\trp(\theta')\,n>0.
\end{equation}
Thus, we can drop the boundary terms and $I_{33}$ can be estimated as
\begin{equation*}
    I_{33}\geq\frac{1}{2}\int_{\mathcal{M}\times\mathcal{N}}\tau\varphi\,D_{t}(\theta')\,\trp(|z|^{2})+\frac{\Delta t}{2}\int_{\mathcal{M}\times\mathcal{N}^{\ast}}\tau\varphi\trm(\theta')|D_{t}z|^{2}.
\end{equation*}    
Now, let us focus on the integral with the term $|\trp z|^{2}$. Using that $|\trp z|^{2}\leq C|\trm z|^{2}+C(\Delta t)^{2}|D_{t}z|^{2}$ and Theorem \ref{theo:fully:weight:estimates} we obtain
\begin{equation*}
    \begin{split}
        I_{33}\geq &-\int_{\mathcal{M}\times\mathcal{N}}\trm(\mu_{33})|\trm z|^{2}-\int_{\mathcal{M}\times\mathcal{N}}(\Delta t)^{2}\trm(\mu_{23})|D_{t}z|^{2}-\int_{\mathcal{M}\times\mathcal{N}}\trm(\gamma_{23})|D_{t}z|^{2},
    \end{split}
\end{equation*}
with $\mu_{33}:=(\tau T^{2}\theta^{3}+\frac{\tau \Delta t}{\delta^{4}T^{5}})\mathcal{O}_{\lambda}(1)$ and $\gamma_{33}:=\Delta t \tau T\theta^{2}\mathcal{O}_{\lambda}(1)$. We can shift the integral with the term $|\trm z|^{2}$.
%%%%%%%%%%%%

%%%%%%%%%%%%%%%%%%%%%
\bibliographystyle{abbrv}
\bibliography{references}
\end{document}
